%% ==================================================
%% Chapter: Dividing a region into dyadic cells
%% ==================================================
\chapter{Dividing a region into dyadic cells}%
\label{ch:peps_area_law_geometry}

The last stage of the proof of the area law \cite[Section~11]{OpenAI2026AreaLaw}
divides the region, apart from a set of sites proportional to the size of its
boundary, into two ordered families of smaller regions. Each region is
separated by a collar from the complement and from the earlier regions of its
own family. This chapter and the next three construct the division. This
chapter covers the dyadic layers of square cells near the boundary, the fine
cells and sparse belts inside them, the primary regions, and the finite family
of marked points of the belt cells.

\section{Sparse coordinate strips}
\label{sec:area_law_sparse_strips}

\begin{definition}[Coordinate residue strips]
    \label{def:area_law_residue_strips}
    Let $F\subset\mathbb Z^2$ be finite and let $m\ge1$ be an integer.
    For residues $a,b\in\{0,\ldots,m-1\}$, define
    \begin{align}
        B_{a,b}(F) & =\{z\in F:z_1\equiv a\pmod m\ \text{or}\ z_2\equiv b\pmod m\}.
        \label{eq:area_law_residue_strips}
    \end{align}
\end{definition}

\begin{theorem}[Sparse choice of residue strips]
    \label{thm:area_law_sparse_belt_shift}
    \lean{TNLean.PEPS.AreaLaw.Geometry.exists_sparse_belt_shift}
    \leanok
    \uses{def:area_law_residue_strips}
    For every finite $F\subset\mathbb Z^2$ and integer $m\ge1$, there are
    residues $a,b\in\{0,\ldots,m-1\}$ such that
    \begin{align}
        m|B_{a,b}(F)| & \le2|F|.
        \label{eq:area_law_sparse_belt_shift}
    \end{align}
    This is the finite averaging step in the sparse-strip construction of
    \cite[Section~11]{OpenAI2026AreaLaw}.
    % Source revision: openai/math@adc7f1241b42e322a6451854ab7e4b4c146bf78a,
    % build/sections/10-geometry.tex, lines 210--237, geometry:belt-count.
    % This auxiliary theorem does not establish the layer-cell count or the
    % final dyadic decay estimate in geometry:belt-count.
\end{theorem}
\begin{proof}\leanok
    For any map $f:F\to\{0,\ldots,m-1\}$, its $m$ fibres partition $F$.
    One fibre therefore has at most $\lfloor|F|/m\rfloor$ elements.
    Apply this to the two coordinate residue maps and choose their
    residues independently. Each selected strip contains at most
    $\lfloor|F|/m\rfloor$ cells, so their union satisfies
    $|B_{a,b}(F)|\le2\lfloor|F|/m\rfloor$. Multiplying by $m$ and using
    $m\lfloor|F|/m\rfloor\le|F|$ gives
    (\ref{eq:area_law_sparse_belt_shift}). The argument includes $F=\varnothing$
    and $m=1$.
\end{proof}

\section{Finite cell counts}
\label{sec:area_law_finite_cell_counts}

\begin{definition}[Boundary endpoints and dyadic children]
    \label{def:area_law_cell_counts}
    \lean{
        TNLean.PEPS.AreaLaw.Geometry.boundaryEndpoints,
        TNLean.PEPS.AreaLaw.Geometry.dyadicChildren
    }
    \leanok
    \uses{def:area_law_hamiltonian}
    For a finite induced domain $\Lambda$ and a cut $A\subseteq\Lambda$, the
    boundary endpoints are the integer endpoints of the crossing edges,
    \begin{align}
        Z & =\bigcup_{\{x,y\}\in\partial_\Lambda A}\{x,y\}\subset\mathbb Z^2.
        \label{eq:area_law_boundary_endpoints}
    \end{align}
    For a finite set $P\subset\mathbb Z^2$ of parent indices, its dyadic child
    indices are
    \begin{align}
        \operatorname{Ch}(P) & =
        \{2z+\varepsilon:z\in P,\ \varepsilon\in\{0,1\}^2\}.
        \label{eq:area_law_dyadic_children}
    \end{align}
    These are the four children of each parent in a translated dyadic mesh.
\end{definition}

\begin{theorem}[Finite cell counts]
    \label{thm:area_law_cell_counts}
    \lean{
        TNLean.PEPS.AreaLaw.Geometry.card_ambientDilation_le,
        TNLean.PEPS.AreaLaw.Geometry.card_boundaryEndpoints_le,
        TNLean.PEPS.AreaLaw.Geometry.card_dyadicChildren
    }
    \leanok
    \uses{def:area_law_cell_counts,def:area_law_template,def:area_law_hamiltonian}
    For every finite $T\subset\mathbb Z^2$, integer $r\ge0$, finite parent
    set $P$, and finite induced domain with cut $A$, one has
    \begin{align}
        |(T+[-r,r]^2)\cap\mathbb Z^2| & \le(2r+1)^2|T|,
        \label{eq:area_law_dilation_count}                        \\
        |Z|                           & \le2|\partial_\Lambda A|,
        \label{eq:area_law_endpoint_count}                        \\
        |\operatorname{Ch}(P)|        & =4|P|.
        \label{eq:area_law_children_count}
    \end{align}
    These cardinality estimates are ingredients in the layer count of
    \cite[Section~11]{OpenAI2026AreaLaw}; no connectedness assumption is needed.
\end{theorem}
\begin{proof}\leanok
    Each point of $T$ contributes an integer square with $(2r+1)^2$ sites;
    counting their union gives the first bound. Each crossing edge contributes
    at most two endpoints, giving the second. In each coordinate the remainder
    modulo two identifies the child label, and division by two then identifies
    the parent index. Thus distinct parent and child-label pairs give distinct
    indices, proving the last equality. All three assertions include empty sets.
\end{proof}

\section{Translated dyadic neighborhoods}
\label{sec:area_law_dyadic_neighborhoods}

\begin{definition}[Half-open dyadic cells]
    \label{def:area_law_dyadic_cells}
    \lean{TNLean.PEPS.AreaLaw.Geometry.dyadicCell,
        TNLean.PEPS.AreaLaw.Geometry.dyadicCellIndex,
        TNLean.PEPS.AreaLaw.Geometry.dyadicParent}
    \leanok
    Fix an origin $o\in\mathbb R^2$. For every integer $k\ge0$ and index
    $z\in\mathbb Z^2$, define the half-open cell and its index map by
    \begin{align}
        C_{k,z} & =o+2^k\bigl(z+[0,1)^2\bigr),                         \\
        c_k(x)  & =\left(\left\lfloor\frac{x_1-o_1}{2^k}\right\rfloor,
        \left\lfloor\frac{x_2-o_2}{2^k}\right\rfloor\right),           \\
        p(z)    & =\left(\left\lfloor\frac{z_1}{2}\right\rfloor,
        \left\lfloor\frac{z_2}{2}\right\rfloor\right).
    \end{align}
    The map $p$ gives the index of a cell's parent. The arbitrary origin
    includes the translated origin used in \cite[Section~11]{OpenAI2026AreaLaw}.
\end{definition}

\begin{theorem}[Cell indices and parents]
    \label{thm:area_law_dyadic_indices}
    \lean{TNLean.PEPS.AreaLaw.Geometry.mem_dyadicCell_iff,
        TNLean.PEPS.AreaLaw.Geometry.dyadicCellIndex_succ,
        TNLean.PEPS.AreaLaw.Geometry.dyadicParent_abs_sub_le,
        TNLean.PEPS.AreaLaw.Geometry.dyadicParent_mem_ambientDilation_image,
        TNLean.PEPS.AreaLaw.Geometry.mem_dyadicChildren_iff}
    \leanok
    \uses{def:area_law_dyadic_cells,def:area_law_cell_counts,def:area_law_template}
    Every point lies in the unique cell indexed by $c_k(x)$, and passing to
    parents agrees with increasing the scale:
    \begin{align}
        x\in C_{k,z} & \quad\Longleftrightarrow\quad c_k(x)=z, \\
        c_{k+1}(x)   & =p(c_k(x)).
    \end{align}
    For every integer $C_0\ge0$, two integer indices at sup distance at most
    $C_0$ have parents at sup distance at most $C_0$. Consequently $p$ maps
    the integer $C_0$-neighborhood of a finite index set into the
    $C_0$-neighborhood of its parent image. An index is a child
    of a member of $P$ precisely when its parent belongs to $P$.
\end{theorem}
\begin{proof}\leanok
    The defining inequalities of the floor function give the cell membership
    criterion. The identity $\lfloor t/2\rfloor=\lfloor\lfloor t\rfloor/2\rfloor$
    gives the parent formula. Integer division by two preserves the stated
    coordinate bounds, including negative indices. Finally the remainder modulo
    two is either zero or one, identifying the unique child in each coordinate.
\end{proof}

\begin{definition}[Occupied indices, neighborhoods and layers]
    \label{def:area_law_dyadic_layers}
    \lean{TNLean.PEPS.AreaLaw.Geometry.occupiedCellIndices,
        TNLean.PEPS.AreaLaw.Geometry.dyadicNeighborhood,
        TNLean.PEPS.AreaLaw.Geometry.dyadicLayerIndices,
        TNLean.PEPS.AreaLaw.Geometry.dyadicLayer}
    \leanok
    \uses{def:area_law_dyadic_cells,def:area_law_cell_counts,def:area_law_template}
    Let $Z\subset\mathbb Z^2$ be finite and let $C_0\ge0$ be an integer.
    For each integer $k\ge0$, set
    \begin{align}
        \mathcal Z_k & =\{c_k(z):z\in Z\},                          \\
        P_k          & =(\mathcal Z_k+[-C_0,C_0]^2)\cap\mathbb Z^2, \\
        N_k          & =\bigcup_{z\in P_k}C_{k,z},                  \\
        D_k          & =N_{k+1}\setminus N_k,                       \\
        I_k          & =\operatorname{Ch}(P_{k+1})\setminus P_k.
    \end{align}
    These are actual subsets of the plane and finite sets of indices; no
    nesting or distance bound is assumed in their definitions.
    The picture below shows one endpoint with $C_0=1$: the layer $D_k$ is
    the ring of scale-$k$ cells that lie in $N_{k+1}$ but not in $N_k$.
    % Source: 10-geometry.tex, lines 148--200 (N_k, D_k = N_{k+1} minus N_k);
    % no source figure. Specialization drawn: Z is one endpoint and C_0 = 1.
    % Ink-to-index: every dot is one half-open scale-k cell C_{k,z}; the red dot
    % is the occupied index c_k(z) in Z_k. The pink enclosure is the 3-by-3
    % block P_k, so its union is N_k. Each dashed square is one parent cell
    % C_{k+1,p} of P_{k+1}, i.e. its four scale-k children; the nine of them
    % make N_{k+1}, marked by the bracket. The tinted ring is the layer D_k,
    % the union of the cells in I_k = Ch(P_{k+1}) minus P_k. N_k is not aligned
    % with the parent squares, which is why D_k cuts through three of them.
    % No strings, legs or contractions: this is a geometry picture with empty
    % boundary signature. Not to scale; the plane continues past the frame.
    \begin{center}
        \begin{tenkz}[size=s, rows={ket,ket,ket,ket,ket,ket,ket,ket}, cols=8, bonds=none]
            \tndeclare{species}{endpoint}{hue=source:red}
            \tn[at=(5,5), skin=dot, species=endpoint]{}
            \tnmark[form=enclosure, species=secondary, tint]{(2,2) .. (7,7) - (4,4) .. (6,6)}{}
            \tnmark[form=enclosure, species=selected, tint]{(4,4) .. (6,6)}{}
            \tnmark[form=enclosure, species=complement]{(2,2) .. (3,3)}{}
            \tnmark[form=enclosure, species=complement]{(2,4) .. (3,5)}{}
            \tnmark[form=enclosure, species=complement]{(2,6) .. (3,7)}{}
            \tnmark[form=enclosure, species=complement]{(4,2) .. (5,3)}{}
            \tnmark[form=enclosure, species=complement]{(4,4) .. (5,5)}{}
            \tnmark[form=enclosure, species=complement]{(4,6) .. (5,7)}{}
            \tnmark[form=enclosure, species=complement]{(6,2) .. (7,3)}{}
            \tnmark[form=enclosure, species=complement]{(6,4) .. (7,5)}{}
            \tnmark[form=enclosure, species=complement]{(6,6) .. (7,7)}{}
            \tnmark[form=label, label pos=e]{midway (2,2) and (3,2)}{$D_k$}
            \tnmark[form=label, label pos=e]{midway (4,4) and (5,4)}{$N_k$}
            \tnmark[form=bracket, label pos=e]{(2,2) .. (7,7)}{$N_{k+1}$}
        \end{tenkz}
    \end{center}
\end{definition}

\begin{theorem}[Nested neighborhoods and layer cells]
    \label{thm:area_law_dyadic_nesting}
    \lean{TNLean.PEPS.AreaLaw.Geometry.occupiedCellIndices_succ,
        TNLean.PEPS.AreaLaw.Geometry.mem_dyadicNeighborhood_iff,
        TNLean.PEPS.AreaLaw.Geometry.dyadicNeighborhood_subset_succ,
        TNLean.PEPS.AreaLaw.Geometry.dyadicLayer_eq_iUnion}
    \leanok
    \uses{def:area_law_dyadic_layers,thm:area_law_dyadic_indices}
    For every origin, finite $Z$, integer $C_0\ge0$, and integer $k\ge0$,
    \begin{align}
        \mathcal Z_{k+1} & =p(\mathcal Z_k),                            \\
        x\in N_k         & \quad\Longleftrightarrow\quad c_k(x)\in P_k, \\
        N_k              & \subseteq N_{k+1},                           \\
        D_k              & =\bigcup_{z\in I_k}C_{k,z}.
    \end{align}
    The statements include the empty set of endpoints and radius zero.
\end{theorem}
\begin{proof}\leanok
    Apply the parent identity to every point of $Z$. The cell membership
    criterion identifies membership in $N_k$ with membership of its index in
    $P_k$. The parent preserves the integer radius, proving nesting.
    A point belongs to $N_{k+1}$ precisely when its scale-$k$ index is a child
    of an index in $P_{k+1}$. Removing $N_k$ therefore removes precisely the
    indices in $P_k$, which gives the last equality.
\end{proof}

\begin{theorem}[Uniform number of layer cells]
    \label{thm:area_law_dyadic_layer_count}
    \lean{TNLean.PEPS.AreaLaw.Geometry.card_dyadicLayerIndices_le,
        TNLean.PEPS.AreaLaw.Geometry.card_dyadicLayerIndices_boundary_le}
    \leanok
    \uses{thm:area_law_cell_counts,thm:area_law_dyadic_nesting}
    For every origin, finite $Z$, integer $C_0\ge0$, and integer $k\ge0$,
    \begin{align}
        |I_k| & \le4(2C_0+1)^2|Z|.
    \end{align}
    When $Z$ is the set of endpoints of the crossing edges of a cut in a
    finite induced domain, with $b=|\partial_\Lambda A|$, this gives
    \begin{align}
        |I_k| & \le8(2C_0+1)^2b.
    \end{align}
    The constant depends only on the chosen radius, independently of the
    origin, scale, domain and cut. These are the layer-cell estimates in
    \cite[Section~11]{OpenAI2026AreaLaw}.
\end{theorem}
\begin{proof}\leanok
    The occupied-index image has at most $|Z|$ elements. Its integer
    $C_0$-neighborhood has at most $(2C_0+1)^2|Z|$ elements, and each parent
    has exactly four children. Taking a set difference can only decrease the
    count. Each crossing edge contributes at most two endpoints, giving the
    second bound.
\end{proof}


\begin{theorem}[Exhaustion of the plane]
    \label{thm:area_law_dyadic_exhaustion}
    \lean{TNLean.PEPS.AreaLaw.Geometry.exists_mem_dyadicNeighborhood,
        TNLean.PEPS.AreaLaw.Geometry.iUnion_dyadicNeighborhood_eq_univ}
    \leanok
    \uses{def:area_law_dyadic_layers,thm:area_law_dyadic_indices}
    For every origin, nonempty finite endpoint set $Z$, and integer $C_0\ge2$,
    every point belongs to some $N_k$, with $k\ge0$. Consequently,
    \begin{align}
        \bigcup_{k\ge0}N_k & =\mathbb R^2.
    \end{align}
    This is the exhaustion assertion in \cite[Section~11]{OpenAI2026AreaLaw}.
\end{theorem}
\begin{proof}\leanok
    Choose $z\in Z$. For any $x\in\mathbb R^2$, take $2^k$ larger than the
    sum of the four absolute values of the coordinates of $x-o$ and $z-o$.
    Each translated coordinate divided by $2^k$ then lies in $(-1,1)$, so
    its floor is either $-1$ or $0$. The indices $c_k(x)$ and $c_k(z)$
    therefore differ by at most one in each coordinate. Since $c_k(z)$ is
    occupied and $C_0\ge2$, the index $c_k(x)$ belongs to $P_k$. Thus
    $x\in N_k$, as required.
\end{proof}


\begin{theorem}[Closures of cells and layers]
    \label{thm:area_law_dyadic_closures}
    \lean{TNLean.PEPS.AreaLaw.Geometry.closure_dyadicCell,
        TNLean.PEPS.AreaLaw.Geometry.convex_closure_dyadicCell,
        TNLean.PEPS.AreaLaw.Geometry.closure_dyadicLayer_eq_iUnion}
    \leanok
    \uses{thm:area_law_dyadic_nesting}
    For every origin, nonnegative scale and cell index, the closure of a
    half-open dyadic cell is the corresponding closed rectangle. For every
    finite endpoint set and nonnegative integer radius, the actual layer satisfies
    \begin{align}
        \overline{C_{k,z}} & =o+2^k\bigl(z+[0,1]^2\bigr),           \\
        \overline{D_k}     & =\bigcup_{z\in I_k}\overline{C_{k,z}}.
    \end{align}
    In particular, the closure of every dyadic cell is convex.
    These formulas give the closed regions used in the distance estimates of
    \cite[Section~11]{OpenAI2026AreaLaw}.
\end{theorem}
\begin{proof}\leanok
    The side length $2^k$ is positive, so the closure of each half-open coordinate
    interval is its closed interval. Closure commutes with finite products.
    The exact layer union is finite, and closure commutes with finite unions,
    which gives the second formula. Each closed coordinate interval is convex,
    and a product of convex sets is convex.
\end{proof}

\section{The translated origin and integer lattice points}
\label{sec:area_law_dyadic_origin}

The geometric construction in \cite[Section~11]{OpenAI2026AreaLaw}
uses a translated origin whose coordinates, sum and difference are
nonintegral.

% Source: 10-geometry.tex, prop:two-families, lines 150--154 and 545--559,
% revision adc7f1241b42e322a6451854ab7e4b4c146bf78a.
% Actual edge containment in these four line families is a separate result.

\begin{definition}[The prescribed origin]
    \label{def:area_law_dyadic_origin}
    \lean{TNLean.PEPS.AreaLaw.Geometry.dyadicOrigin}
    \leanok
    Set $o=(\sqrt 2,\sqrt 3)\in\mathbb R^2$.
\end{definition}

\begin{lemma}[Nonintegrality of the origin coordinates]
    \label{thm:area_law_dyadic_origin_nonintegral}
    \lean{TNLean.PEPS.AreaLaw.Geometry.dyadicOrigin_nonintegral}
    \leanok
    \uses{def:area_law_dyadic_origin}
    None of $o_1$, $o_2$, $o_1+o_2$ and $o_1-o_2$ is an integer.
\end{lemma}
\begin{proof}\leanok
    \uses{def:area_law_dyadic_origin}
    Squaring positive rational bounds gives
    $4/3<\sqrt 2<2$ and $5/3<\sqrt 3<2$; also
    $\sqrt 2<\sqrt 3$. Thus $o_1,o_2\in(1,2)$,
    $o_1+o_2\in(3,4)$ and $o_1-o_2\in(-1,0)$,
    and no integer lies strictly between consecutive integers.
\end{proof}

\begin{lemma}[Integer-translated supporting lines avoid the lattice]
    \label{thm:area_law_dyadic_origin_lines}
    \lean{TNLean.PEPS.AreaLaw.Geometry.dyadicOrigin_supporting_lines_avoid_lattice}
    \leanok
    \uses{def:area_law_dyadic_origin}
    For every $m\in\mathbb Z$, none of the four lines
    \begin{align}
        x_1     & =o_1+m,     \\
        x_2     & =o_2+m,     \\
        x_1+x_2 & =o_1+o_2+m, \\
        x_1-x_2 & =o_1-o_2+m
    \end{align}
    contains a point of $\mathbb Z^2$.
\end{lemma}
\begin{proof}\leanok
    \uses{thm:area_law_dyadic_origin_nonintegral}
    If $x\in\mathbb Z^2$ satisfies one of these equations, subtracting
    $m$ expresses $o_1$, $o_2$, $o_1+o_2$ or $o_1-o_2$ as an integer,
    which is impossible since none of them is an integer.
\end{proof}

\section{Distances between dyadic layers}
\label{sec:area_law_layer_distance}

Retain the cells, neighborhoods and layers of
Definition~\ref{def:area_law_dyadic_layers}. The origin $o\in\mathbb R^2$ is
arbitrary, $Z\subset\mathbb Z^2$ is finite, and $C_0\ge0$ is an integer.
All scale indices are nonnegative integers. Distances in the plane are taken
in the sup metric, so $d_\infty(x,y)=\max\{|x_1-y_1|,|x_2-y_2|\}$.

\begin{lemma}[Distances between closed cells]
    \label{thm:area_law_dyadic_cell_distance}
    \lean{TNLean.PEPS.AreaLaw.Geometry.dyadicCell_dist_le,
        TNLean.PEPS.AreaLaw.Geometry.dyadicCell_dist_lower}
    \leanok
    \uses{def:area_law_dyadic_cells}
    Let $u,v\in\mathbb Z^2$, $x\in\overline{C_{k,u}}$, and
    $y\in\overline{C_{k,v}}$. If $|u_i-v_i|\le C_0$ for both coordinates
    $i\in\{1,2\}$, then
    \begin{align}
        d_\infty(x,y) & \le (C_0+1)2^k.
    \end{align}
    If $|u_i-v_i|>C_0$ in at least one coordinate, then
    \begin{align}
        C_0 2^k & \le d_\infty(x,y).
    \end{align}
\end{lemma}
\begin{proof}\leanok
    \uses{thm:area_law_dyadic_closures}
    Each coordinate of a closed cell lies in its closed interval of length
    $2^k$. Indices differing by at most $C_0$ therefore give coordinate
    distances at most $(C_0+1)2^k$. If two integer indices differ by more
    than $C_0$, they differ by at least $C_0+1$; the corresponding closed
    intervals are separated by at least $C_0 2^k$. Taking the maximum of
    the two coordinate distances proves both bounds.
\end{proof}

% Source: geometry:layer-distance in Section 11.
\begin{lemma}[Distance from a layer to the endpoints]
    \label{thm:area_law_dyadic_layer_distance}
    \lean{TNLean.PEPS.AreaLaw.Geometry.dyadicLayer_dist_lower,
        TNLean.PEPS.AreaLaw.Geometry.dyadicLayer_exists_dist_le,
        TNLean.PEPS.AreaLaw.Geometry.dyadicLayer_infDist_bounds}
    \leanok
    \uses{def:area_law_dyadic_layers}
    For every $x\in\overline{D_k}$ and every $z\in Z$,
    \begin{align}
        C_0 2^k & \le d_\infty(x,z).
    \end{align}
    For every $x\in\overline{D_k}$ there is a point $z_*\in Z$ such that
    \begin{align}
        d_\infty(x,z_*) & \le 2(C_0+1)2^k.
    \end{align}
    Writing $d_\infty(x,Z)=\inf_{z\in Z}d_\infty(x,z)$, one obtains
    \begin{align}
        C_0 2^k & \le d_\infty(x,Z) \le 2(C_0+1)2^k.
    \end{align}
    This is the layer-distance estimate in
    \cite[Section~11]{OpenAI2026AreaLaw}.
\end{lemma}
\begin{proof}\leanok
    \uses{thm:area_law_dyadic_closures,thm:area_law_dyadic_indices,
        thm:area_law_dyadic_cell_distance}
    The closure of $D_k$ is a finite union of closed scale-$k$ cells.
    Every layer index lies outside $P_k$ and so differs from every occupied
    index in at least one coordinate by more than $C_0$. Applying the
    closed-cell lower bound to the cell containing $z$ gives the first
    inequality.

    Since $\overline{D_k}\subseteq\overline{N_{k+1}}$, the point $x$ lies in a
    closed scale-$(k+1)$ cell whose index is within $C_0$ of an occupied index.
    Choosing $z_*\in Z$ in that occupied cell and applying the closed-cell
    upper bound at scale $k+1$ gives the second inequality.
\end{proof}

% Source: geometry:nonadjacent in Section 11.
\begin{lemma}[Separation of nonadjacent layers]
    \label{thm:area_law_dyadic_nonadjacent}
    \lean{TNLean.PEPS.AreaLaw.Geometry.dyadicLayer_dist_separation,
        TNLean.PEPS.AreaLaw.Geometry.dyadicLayer_dist_nonadjacent}
    \leanok
    \uses{def:area_law_dyadic_layers}
    For every pair of scales $k,h$, every $x\in\overline{D_k}$, and every
    $y\in\overline{D_h}$,
    \begin{align}
        C_0 2^h-2(C_0+1)2^k & \le d_\infty(x,y).
    \end{align}
    If $h\ge k+2$, then
    \begin{align}
        \frac{C_0-1}{2}2^h & \le d_\infty(x,y).
    \end{align}
    This is the nonadjacent-layer estimate in
    \cite[Section~11]{OpenAI2026AreaLaw}.
\end{lemma}
\begin{proof}\leanok
    \uses{thm:area_law_dyadic_layer_distance}
    Choose $z\in Z$ with $d_\infty(x,z)\le2(C_0+1)2^k$. Since
    $d_\infty(y,z)\ge C_0 2^h$, the triangle inequality gives the first bound.
    When $h\ge k+2$, $4\cdot2^k\le2^h$, so
    \begin{align}
        C_0 2^h-2(C_0+1)2^k & \ge \frac{C_0-1}{2}2^h.
    \end{align}
\end{proof}

\section{The dummy neighborhood meets only the initial layer}
\label{sec:area_law_dummy_contacts}

The dummy tile is the initial neighborhood $N_{k_0}$. Its distance
bound distinguishes it from the later layers; see
\cite[Section~11]{OpenAI2026AreaLaw}. The conclusions below concern
closures, with distances measured in the sup metric.

% Source: 10-geometry.tex, lines 160--191 and 299--306.

\begin{lemma}[Distance from the neighborhood to an endpoint]
    \label{thm:area_law_neighborhood_endpoint_distance}
    \lean{TNLean.PEPS.AreaLaw.Geometry.dyadicNeighborhood_exists_dist_le}
    \leanok
    \uses{def:area_law_dyadic_layers}
    For any origin, finite endpoint set $Z$, integer width $C_0\ge 0$,
    and layer index $k$, every $x\in\overline{N_k}$ admits
    an endpoint $z\in Z$ with
    \begin{align}
        d_\infty(x,z)\le(C_0+1)2^k.
    \end{align}
\end{lemma}
\begin{proof}\leanok
    \uses{thm:area_law_dyadic_closures,thm:area_law_dyadic_indices,
        thm:area_law_dyadic_cell_distance}
    The point lies in a closed scale-$k$ cell whose index is within
    $C_0$ of an occupied index in each coordinate. Choose an endpoint
    in that occupied cell and apply the closed-cell distance bound.
\end{proof}

\begin{lemma}[Separation from subsequent layers]
    \label{thm:area_law_dummy_later_layer_separation}
    \lean{TNLean.PEPS.AreaLaw.Geometry.dyadicNeighborhood_dist_later_layer}
    \leanok
    \uses{def:area_law_dyadic_layers}
    Under the same hypotheses, let $k\ge k_0+1$,
    $x\in\overline{N_{k_0}}$, and $y\in\overline{D_k}$. Then
    \begin{align}
        (C_0-1)2^{k_0}\le d_\infty(x,y).
    \end{align}
    In particular, this is a positive separation when $C_0\ge2$.
\end{lemma}
\begin{proof}\leanok
    \uses{thm:area_law_neighborhood_endpoint_distance,
        thm:area_law_dyadic_layer_distance}
    Choose $z\in Z$ with
    $d_\infty(x,z)\le(C_0+1)2^{k_0}$. The layer bound gives
    $C_0 2^k\le d_\infty(y,z)$, while $2^k\ge2\cdot2^{k_0}$.
    The triangle inequality now yields
    \begin{align}
        d_\infty(x,y) & \ge C_0 2^k-(C_0+1)2^{k_0}
        \ge(C_0-1)2^{k_0}.
    \end{align}
\end{proof}

\begin{lemma}[Only the initial layer can meet the dummy]
    \label{thm:area_law_dummy_contact_index}
    \lean{TNLean.PEPS.AreaLaw.Geometry.dyadicNeighborhood_contact_layer_eq}
    \leanok
    \uses{def:area_law_dyadic_layers}
    Suppose that $C_0\ge2$ and $k\ge k_0$. If
    $\overline{N_{k_0}}\cap\overline{D_k}$ is nonempty, then $k=k_0$.
\end{lemma}
\begin{proof}\leanok
    \uses{thm:area_law_dummy_later_layer_separation}
    If $k\ne k_0$, then $k\ge k_0+1$. A shared point has distance
    zero from itself, contradicting the positive separation.
\end{proof}

\section{Fine cells and sparse belts}
\label{sec:area_law_fine_belts}

The origin $o\in\mathbb R^2$ remains arbitrary. All scale indices and
refinement depths are nonnegative integers. Let $Z\subset\mathbb Z^2$ be
finite, and let $C_0\ge0$ be an integer. Retain the layer $D_k$ and its
index set $I_k$ from Definition~\ref{def:area_law_dyadic_layers}.

\begin{definition}[Dyadic refinement]
    \label{def:area_law_dyadic_refinement}
    \lean{TNLean.PEPS.AreaLaw.Geometry.dyadicAncestor,
        TNLean.PEPS.AreaLaw.Geometry.dyadicRefinement}
    \leanok
    For every refinement depth $d$ and fine index $u\in\mathbb Z^2$, define
    its coarse ancestor by
    \begin{align}
        A_d(u) & =\left(\left\lfloor\frac{u_1}{2^d}\right\rfloor,
        \left\lfloor\frac{u_2}{2^d}\right\rfloor\right).
    \end{align}
    For every finite set $P\subset\mathbb Z^2$, define its refinement by
    \begin{align}
        R_d(P) & =\{2^d z+a:z\in P,\ a\in\{0,\ldots,2^d-1\}^2\}.
    \end{align}
\end{definition}

\begin{theorem}[Ancestors and refinement counts]
    \label{thm:area_law_dyadic_refinement}
    \lean{TNLean.PEPS.AreaLaw.Geometry.mem_dyadicRefinement_iff,
        TNLean.PEPS.AreaLaw.Geometry.card_dyadicRefinement,
        TNLean.PEPS.AreaLaw.Geometry.dyadicCellIndex_of_le}
    \leanok
    \uses{def:area_law_dyadic_refinement,def:area_law_dyadic_cells}
    For every finite $P$, depth $d$, and index $u\in\mathbb Z^2$,
    \begin{align}
        u\in R_d(P) & \quad\Longleftrightarrow\quad A_d(u)\in P, \\
        |R_d(P)|    & =4^d|P|.
    \end{align}
    For every pair of scales $\ell\le k$ and every $x\in\mathbb R^2$,
    \begin{align}
        A_{k-\ell}(c_\ell(x)) & =c_k(x).
    \end{align}
\end{theorem}
\begin{proof}\leanok
    Euclidean division by $2^d$ gives a unique integer quotient and a
    remainder in $\{0,\ldots,2^d-1\}$ in each coordinate. This proves the
    membership criterion and the cardinality formula, including negative
    indices. The identity $\lfloor t/n\rfloor=\lfloor\lfloor t\rfloor/n\rfloor$
    for a positive integer $n$ gives the point-index identity.
\end{proof}

\begin{definition}[Fine indices of a layer]
    \label{def:area_law_fine_layer_indices}
    \lean{TNLean.PEPS.AreaLaw.Geometry.fineLayerIndices}
    \leanok
    \uses{def:area_law_dyadic_refinement,def:area_law_dyadic_layers}
    For every fine scale $\ell\le k$, set
    \begin{align}
        F_{k,\ell} & =R_{k-\ell}(I_k).
    \end{align}
\end{definition}

\begin{theorem}[Exact fine-cell layer union]
    \label{thm:area_law_fine_layer_count}
    \lean{TNLean.PEPS.AreaLaw.Geometry.dyadicLayer_eq_iUnion_fine,
        TNLean.PEPS.AreaLaw.Geometry.card_fineLayerIndices,
        TNLean.PEPS.AreaLaw.Geometry.card_fineLayerIndices_le,
        TNLean.PEPS.AreaLaw.Geometry.card_fineLayerIndices_boundary_le}
    \leanok
    \uses{def:area_law_fine_layer_indices,def:area_law_dyadic_cells}
    For every $\ell\le k$, the layer and its fine-cell count satisfy
    \begin{align}
        D_k          & =\bigcup_{u\in F_{k,\ell}} C_{\ell,u}, \\
        |F_{k,\ell}| & =4^{k-\ell}|I_k|,                      \\
        |F_{k,\ell}| & \le4(2C_0+1)^2|Z|\,4^{k-\ell}.
    \end{align}
    If $Z$ is the set of endpoints of the crossing edges of a cut in a
    finite induced domain and $b$ is the number of those edges, then
    \begin{align}
        |F_{k,\ell}| & \le8(2C_0+1)^2b\,4^{k-\ell}.
    \end{align}
    Because the half-open fine cells are disjoint and nonempty, $F_{k,\ell}$
    is the set of fine indices whose cells meet $D_k$, matching the
    cell count of \cite[Section~11]{OpenAI2026AreaLaw}.
\end{theorem}
\begin{proof}\leanok
    \uses{thm:area_law_dyadic_refinement,thm:area_law_dyadic_indices,
        thm:area_law_dyadic_nesting,thm:area_law_dyadic_layer_count}
    A point's fine index belongs to $R_{k-\ell}(I_k)$ if and only if its
    coarse index belongs to $I_k$. The cell-membership criterion therefore
    proves the decomposition of $D_k$. Multiplying the coarse-cell count by
    $4^{k-\ell}$ and applying the endpoint and crossing-edge bounds gives
    the cardinality estimates.
\end{proof}

\begin{definition}[Fine scale and belt pitch]
    \label{def:area_law_fine_belt_scales}
    \lean{TNLean.PEPS.AreaLaw.Geometry.fineScaleIndex,
        TNLean.PEPS.AreaLaw.Geometry.pitchScaleIndex}
    \leanok
    \uses{thm:area_law_exponent_gaps}
    With the fixed exponents $\delta_0$ and $\zeta=1-\delta_0/2$, define,
    for every scale $k$,
    \begin{align}
        \ell_k & =\lfloor\zeta k\rfloor,       \\
        p_k    & =\lfloor(1+\delta_0)k\rfloor, \\
        r_k    & =2^k,                         \\
        t_k    & =2^{\ell_k},                  \\
        s_k    & =2^{p_k}.
    \end{align}
    The fine cells have side $t_k$, while $s_k$ is the belt pitch, as in
    \cite[Section~11]{OpenAI2026AreaLaw}.
\end{definition}

\begin{theorem}[Scale order, divisibility and rounding]
    \label{thm:area_law_fine_belt_scales}
    \lean{TNLean.PEPS.AreaLaw.Geometry.fineScaleIndex_le,
        TNLean.PEPS.AreaLaw.Geometry.le_pitchScaleIndex,
        TNLean.PEPS.AreaLaw.Geometry.fineScale_dvd_layerScale,
        TNLean.PEPS.AreaLaw.Geometry.layerScale_dvd_pitchScale,
        TNLean.PEPS.AreaLaw.Geometry.layerScale_exponent_le}
    \leanok
    \uses{def:area_law_fine_belt_scales}
    For every scale $k$,
    \begin{align}
        \ell_k        & \le k\le p_k,              \\
        t_k           & \mid r_k,                  \\
        r_k           & \mid s_k,                  \\
        2k-p_k-\ell_k & \le2-\frac{\delta_0 k}{2}.
    \end{align}
\end{theorem}
\begin{proof}\leanok
    \uses{thm:area_law_exponent_gaps}
    The inequalities $\zeta\le1$ and $\delta_0\ge0$ give $\ell_k\le k\le p_k$,
    which implies the divisibility relations for the corresponding powers of two.
    Since each floor exceeds its argument minus one, summing the two floor
    bounds and substituting $\zeta=1-\delta_0/2$ yields the exponent bound.
\end{proof}

\begin{definition}[Belt cells]
    \label{def:area_law_belt_cells}
    \lean{TNLean.PEPS.AreaLaw.Geometry.beltCellIndices}
    \leanok
    Let $F\subset\mathbb Z^2$ be finite and let $m\ge1$ be an integer.
    For every pair of residues $a,b\in\{0,\ldots,m-1\}$, define
    \begin{align}
        B(F;m,a,b) & =\{u\in F:u_1\equiv a\pmod m\ \text{or}\ u_2\equiv b\pmod m\}.
    \end{align}
    For a fine layer with pitch exponent $p\ge k\ge\ell$, the modulus is
    $m=2^{p-\ell}$, the ratio of pitch to fine-cell side.
\end{definition}

\begin{theorem}[Sparse shift in the fine-cell layer]
    \label{thm:area_law_fine_belt_shift}
    \lean{TNLean.PEPS.AreaLaw.Geometry.exists_fine_belt_shift}
    \leanok
    \uses{def:area_law_belt_cells,def:area_law_fine_layer_indices}
    Let $Z$ be the crossing-edge endpoint set of a finite induced-domain
    cut, with crossing-edge count $b$. For every fine scale $\ell\le k$
    and every pitch exponent $p\ge0$, put $F=F_{k,\ell}$ and
    $m=2^{\max(p-\ell,0)}$.
    There are residues $a,b'\in\{0,\ldots,m-1\}$ such that
    \begin{align}
        m\,|B(F;m,a,b')| & \le16(2C_0+1)^2b\,4^{k-\ell}.
    \end{align}
    When $p\ge\ell$, put $r=2^k$, $t=2^\ell$, and $s=2^p$; then $m=s/t$
    and
    \begin{align}
        |B(F;m,a,b')| & \le16(2C_0+1)^2b\,\frac{r^2}{st}.
    \end{align}
\end{theorem}
\begin{proof}\leanok
    \uses{thm:area_law_sparse_belt_shift,thm:area_law_fine_layer_count}
    Averaging over coordinate residues mod $m$ selects residues $(a,b')$ with
    $m|B(F;m,a,b')|\le2|F|$. The fine-cell count gives the first inequality;
    dividing by $m=2^{p-\ell}$ and using $4^{k-\ell}/2^{p-\ell}=r^2/(st)$
    gives the second.
\end{proof}

\begin{theorem}[Geometric decay of the scale ratio]
    \label{thm:area_law_fine_belt_ratio}
    \lean{TNLean.PEPS.AreaLaw.Geometry.dyadicScale_ratio_le}
    \leanok
    \uses{def:area_law_fine_belt_scales}
    For every scale $k$,
    \begin{align}
        \frac{r_k^2}{s_k t_k} & \le4\,2^{-\delta_0 k/2}.
    \end{align}
\end{theorem}
\begin{proof}\leanok
    \uses{thm:area_law_fine_belt_scales}
    The ratio equals $2^{2k-p_k-\ell_k}$. The exponent bound gives
    $2^{2k-p_k-\ell_k}\le2^{2-\delta_0 k/2}=4\,2^{-\delta_0 k/2}$.
\end{proof}

% Source: geometry:belt-count in Section 11.
\begin{theorem}[Sparse dyadic belts]
    \label{thm:area_law_sparse_dyadic_belts}
    \lean{TNLean.PEPS.AreaLaw.Geometry.exists_sparse_dyadic_belt_shift}
    \leanok
    \uses{def:area_law_belt_cells,def:area_law_fine_layer_indices,
        def:area_law_fine_belt_scales}
    Let $Z$ be the crossing-edge endpoint set of a finite induced-domain
    cut, with crossing-edge count $b$. For every scale $k$, put
    $F=F_{k,\ell_k}$ and $m=2^{p_k-\ell_k}=s_k/t_k$. There are residues
    $a,b'\in\{0,\ldots,m-1\}$ such that
    \begin{align}
        |B(F;m,a,b')| & \le64(2C_0+1)^2b\,2^{-\delta_0 k/2}.
    \end{align}
    This is the sparse-belt estimate in \cite[Section~11]{OpenAI2026AreaLaw}.
    The constant depends only on $C_0$, and the residues are chosen
    independently in each layer.
\end{theorem}
\begin{proof}\leanok
    \uses{thm:area_law_fine_belt_shift,thm:area_law_fine_belt_scales,
        thm:area_law_fine_belt_ratio}
    Apply the belt-shift bound at exponent $\ell=\ell_k$ and pitch
    $p=p_k$, and substitute the scale-ratio bound for $r_k^2/(s_k t_k)$.
\end{proof}

\section{The half-open layer partition}
\label{sec:area_law_layer_partition}

The nested neighborhoods $N_k$ define the half-open layers
$D_k=N_{k+1}\setminus N_k$, which partition $\mathbb R^2\setminus N_{k_0}$
before closures are taken; see \cite[Section~11]{OpenAI2026AreaLaw}.

% Source: 10-geometry.tex, lines 154--181. All disjointness statements
% concern the half-open sets, not their closures.

\begin{lemma}[Fine cells lie in their layer]
    \label{thm:area_law_fine_cell_containment}
    \lean{TNLean.PEPS.AreaLaw.Geometry.dyadicCell_subset_dyadicLayer_of_mem_fineLayerIndices}
    \leanok
    \uses{def:area_law_fine_layer_indices}
    For any origin $o$, finite endpoint set $Z$, and integer $C_0\ge0$, if
    $\ell\le k$ and $z\in F_{k,\ell}$, then $C_{\ell,z}\subseteq D_k$.
\end{lemma}
\begin{proof}\leanok
    \uses{thm:area_law_fine_layer_count}
    The layer $D_k$ is the union of $C_{\ell,w}$ over $w\in F_{k,\ell}$.
\end{proof}

\begin{lemma}[Pairwise disjoint half-open layers]
    \label{thm:area_law_layer_pairwise_disjoint}
    \lean{TNLean.PEPS.AreaLaw.Geometry.dyadicLayer_pairwiseDisjoint}
    \leanok
    \uses{def:area_law_dyadic_layers}
    For any origin $o$, finite endpoint set $Z$, and integer $C_0\ge0$,
    $D_k\cap D_h=\varnothing$ whenever $k\ne h$.
\end{lemma}
\begin{proof}\leanok
    \uses{def:area_law_dyadic_layers,thm:area_law_dyadic_nesting}
    If $k<h$, nesting of the neighborhoods gives
    $D_k\subseteq N_{k+1}\subseteq N_h$, whereas
    $D_h=N_{h+1}\setminus N_h$ is disjoint from $N_h$. The case $h<k$ follows
    by symmetry.
\end{proof}

\begin{lemma}[The dummy region is disjoint from later layers]
    \label{thm:area_law_dummy_layer_disjoint}
    \lean{TNLean.PEPS.AreaLaw.Geometry.dyadicNeighborhood_disjoint_later_layer}
    \leanok
    \uses{def:area_law_dyadic_layers}
    Under the same hypotheses, for every initial layer index $k_0$ and every
    $k\ge k_0$, the neighborhood $N_{k_0}$ is disjoint from $D_k$.
\end{lemma}
\begin{proof}\leanok
    \uses{def:area_law_dyadic_layers,thm:area_law_dyadic_nesting}
    Nesting gives $N_{k_0}\subseteq N_k$, and $D_k=N_{k+1}\setminus N_k$.
\end{proof}

\begin{lemma}[Unique layer outside the dummy region]
    \label{thm:area_law_unique_layer}
    \lean{TNLean.PEPS.AreaLaw.Geometry.exists_unique_layer_of_not_mem_dyadicNeighborhood}
    \leanok
    \uses{def:area_law_dyadic_layers}
    Suppose that $Z$ is nonempty and $C_0\ge2$. For every initial layer
    index $k_0$ and every $x\notin N_{k_0}$, there is a unique
    $k\ge k_0$ such that $x\in D_k$.
\end{lemma}
\begin{proof}\leanok
    \uses{thm:area_law_dyadic_exhaustion,thm:area_law_dyadic_nesting,
        thm:area_law_layer_pairwise_disjoint}
    Exhaustion of $\mathbb R^2$ by the neighborhoods $N_j$ yields a least
    $K\ge0$ with $x\in N_K$. Since $x\notin N_{k_0}$ and the neighborhoods
    $N_j$ are nested, $K>k_0$. Minimality of $K$ gives $x\notin N_{K-1}$,
    so $x\in D_{K-1}$. Uniqueness follows from pairwise disjointness of the
    layers.
\end{proof}

\begin{lemma}[Exact cover by the dummy region and later layers]
    \label{thm:area_law_dummy_layers_cover}
    \lean{TNLean.PEPS.AreaLaw.Geometry.dyadicNeighborhood_union_iUnion_layers_eq_univ}
    \leanok
    \uses{def:area_law_dyadic_layers}
    For nonempty $Z$, integer $C_0\ge2$, and every initial layer index $k_0$,
    \begin{align}
        \mathbb R^2=N_{k_0}\cup\bigcup_{k\ge k_0}D_k.
    \end{align}
\end{lemma}
\begin{proof}\leanok
    \uses{thm:area_law_unique_layer}
    Every point of $\mathbb R^2\setminus N_{k_0}$ lies in $D_k$ for its
    unique layer index $k\ge k_0$.
\end{proof}

\section{Unique fine cells outside the dummy region}
\label{sec:area_law_fine_cell_partition}

The half-open fine cells refine the dyadic layer partition of
\cite[Section~11]{OpenAI2026AreaLaw}. Retain the origin $o$, the finite
endpoint set $Z$, the integer width $C_0$, and the fine exponents $\ell_k$.

% Source: 10-geometry.tex, lines 154--181 and 200--207.
% These statements assign half-open fine cells. They do not assign the
% final primary and belt regions or prove the full two-family proposition.

\begin{lemma}[Unique fine cell outside the initial neighborhood]
    \label{thm:area_law_unique_fine_cell}
    \lean{TNLean.PEPS.AreaLaw.Geometry.exists_unique_fineCell_of_not_mem_dyadicNeighborhood}
    \leanok
    \uses{def:area_law_dyadic_layers,def:area_law_fine_layer_indices,
        def:area_law_fine_belt_scales}
    Suppose that $Z$ is nonempty and $C_0\ge2$. For every initial layer index
    $k_0$ and every $x\notin N_{k_0}$, there is a unique pair $(k,z)$
    with $k\ge k_0$, $z\in F_{k,\ell_k}$, and $x\in C_{\ell_k,z}$.
\end{lemma}
\begin{proof}\leanok
    \uses{thm:area_law_unique_layer,thm:area_law_fine_layer_count,
        thm:area_law_fine_cell_containment,thm:area_law_fine_belt_scales,
        thm:area_law_dyadic_indices}
    The point $x$ belongs to a unique layer $D_k$ with $k\ge k_0$. Since
    $\ell_k\le k$, the decomposition of $D_k$ into scale-$\ell_k$ fine cells
    gives $z\in F_{k,\ell_k}$ with $x\in C_{\ell_k,z}$. Conversely, if
    $h\ge k_0$, $w\in F_{h,\ell_h}$, and $x\in C_{\ell_h,w}$, then
    $C_{\ell_h,w}\subseteq D_h$, so $h=k$, and $w=c_{\ell_k}(x)=z$.
\end{proof}

\begin{lemma}[Exact cover by the dummy region and fine cells]
    \label{thm:area_law_dummy_fine_cells_cover}
    \lean{TNLean.PEPS.AreaLaw.Geometry.dyadicNeighborhood_union_iUnion_fineCells_eq_univ}
    \leanok
    \uses{def:area_law_dyadic_layers,def:area_law_fine_layer_indices,
        def:area_law_fine_belt_scales}
    For nonempty $Z$, integer $C_0\ge2$, and every initial layer index $k_0$,
    \begin{align}
        \mathbb R^2=N_{k_0}\cup\bigcup_{k\ge k_0}
        \bigcup_{z\in F_{k,\ell_k}}C_{\ell_k,z}.
    \end{align}
\end{lemma}
\begin{proof}\leanok
    \uses{thm:area_law_unique_fine_cell}
    Every $x\in\mathbb R^2$ either lies in $N_{k_0}$ or belongs to the fine
    cell $C_{\ell_k,z}$ of its unique pair $(k,z)$ with $k\ge k_0$ and
    $z\in F_{k,\ell_k}$.
\end{proof}

\section{Uniform separation of the dyadic scales}
\label{sec:area_law_scale_separation}

Retain the fixed exponents and scales of
Definition~\ref{def:area_law_fine_belt_scales}. In particular,
$\delta_0=2\cdot10^{-7}$, $\zeta=1-\delta_0/2$, $t_k=2^{\ell_k}$,
$r_k=2^k$, and $s_k=2^{p_k}$. All indices below are nonnegative integers.
The following statements give the uniform scale separation used in
\cite[Section~11]{OpenAI2026AreaLaw}.

\begin{theorem}[Explicit gaps between scale indices]
    \label{thm:area_law_scale_index_gaps}
    \lean{TNLean.PEPS.AreaLaw.Geometry.dyadicScale_index_gaps}
    \leanok
    \uses{def:area_law_fine_belt_scales}
    For every pair of indices $d,k$ satisfying $10^7d\le k$,
    \begin{align}
        \ell_k+d & \le k,   \\
        k+d      & \le p_k.
    \end{align}
\end{theorem}
\begin{proof}\leanok
    \uses{thm:area_law_exponent_gaps}
    The threshold gives $d\le k$ and
    $\zeta k=k-10^{-7}k\le k-d$. Taking the floor proves the first
    inequality. Likewise $(1+\delta_0)k\ge k+2d\ge k+d$, so its floor
    is at least the integer $k+d$.
\end{proof}

\begin{theorem}[Dyadic factors between successive sides]
    \label{thm:area_law_scale_side_ratios}
    \lean{TNLean.PEPS.AreaLaw.Geometry.dyadicScale_side_ratios}
    \leanok
    \uses{def:area_law_fine_belt_scales}
    For every pair of indices $d,k$ satisfying $10^7d\le k$,
    \begin{align}
        2^d t_k & \le r_k, \\
        2^d r_k & \le s_k.
    \end{align}
\end{theorem}
\begin{proof}\leanok
    \uses{thm:area_law_scale_index_gaps}
    Apply the increasing function $n\mapsto2^n$ to the two index
    inequalities and use $2^{a+b}=2^a2^b$.
\end{proof}

\begin{theorem}[Uniform multiplicative separation]
    \label{thm:area_law_uniform_scale_separation}
    \lean{TNLean.PEPS.AreaLaw.Geometry.exists_dyadicScale_side_ratios}
    \leanok
    \uses{def:area_law_fine_belt_scales}
    For every real number $M$, there is a nonnegative integer $K$ such
    that, for every nonnegative integer $k\ge K$,
    \begin{align}
        M t_k & \le r_k, \\
        M r_k & \le s_k.
    \end{align}
    The threshold depends only on $M$ and the fixed exponents,
    independently of the origin, domain, cut, and neighborhood radius.
\end{theorem}
\begin{proof}\leanok
    \uses{thm:area_law_scale_side_ratios}
    Choose a nonnegative integer $d$ with $M<2^d$, using the unboundedness
    of powers of two, and set $K=10^7d$. For every $k\ge K$, the preceding
    theorem gives $2^dt_k\le r_k$ and $2^dr_k\le s_k$. The cell sides are
    positive, so replacing $2^d$ by $M$ proves both inequalities.
\end{proof}

\section{Polynomial factors in the dyadic series}
\label{sec:area_law_polynomial_decay}

With the exponent $\delta_0=2\cdot10^{-7}$, define for each $p\in\mathbb R$
and integer $k\ge0$
\begin{align}
    w_p(k) & =(k+1)^p2^{-\delta_0 k/2}.
\end{align}
The estimates below establish convergence of the series $\sum_{k\ge0}w_p(k)$
used in \cite[Section~11]{OpenAI2026AreaLaw}. Bounding the total number of
repairs by this series further requires bounding the descendants of each
initial mark and relating the initial marks to the sparse belts.

% Source: 10-geometry.tex, lines 232--235 and 668--692,
% geometry:descendant-count and geometry:total-repairs.
% These are auxiliary numerical statements. The construction and count of
% the repairs, and the full statement geometry:total-repairs,
% are treated separately.

\begin{lemma}[Eventual absorption of a polynomial factor]
    \label{thm:area_law_polynomial_absorption}
    \lean{TNLean.PEPS.AreaLaw.Geometry.exists_polynomial_dyadic_absorption}
    \leanok
    \uses{thm:area_law_exponent_gaps}
    For every $p\in\mathbb R$, there exists an integer $K\ge0$, depending only
    on $p$ and $\delta_0$, such that for all integers $k\ge K$,
    \begin{align}
        (k+1)^p2^{-\delta_0 k/2} & \le2^{-\delta_0 k/4}.
    \end{align}
\end{lemma}
\begin{proof}\leanok
    \uses{thm:area_law_exponent_gaps}
    Set $a=\delta_0\log2/4>0$. Since $(k+1)^p\exp(-ak)\to0$ as $k\to\infty$,
    there exists $K\ge0$ such that $(k+1)^p2^{-\delta_0 k/4}\le1$ for all
    $k\ge K$. Multiplying by $2^{-\delta_0 k/4}$ gives the inequality.
\end{proof}

\begin{lemma}[Summability of the polynomially weighted decay]
    \label{thm:area_law_polynomial_summability}
    \lean{TNLean.PEPS.AreaLaw.Geometry.summable_polynomial_dyadic_decay}
    \leanok
    \uses{thm:area_law_exponent_gaps}
    For every $p\in\mathbb R$, the series
    \begin{align}
        \sum_{k=0}^{\infty}(k+1)^p2^{-\delta_0 k/2}
    \end{align}
    converges.
\end{lemma}
\begin{proof}\leanok
    \uses{thm:area_law_polynomial_absorption,thm:area_law_exponent_gaps}
    Set $q=2^{-\delta_0/4}\in(0,1)$. By the preceding theorem,
    $w_p(k)\le q^k$ for all sufficiently large $k$, so convergence follows by
    comparison with the geometric series $\sum_{k\ge0}q^k$.
\end{proof}

\begin{lemma}[Uniform bound on finite scale sums]
    \label{thm:area_law_polynomial_finite_sums}
    \lean{TNLean.PEPS.AreaLaw.Geometry.exists_uniform_polynomial_dyadic_sum_bound}
    \leanok
    \uses{thm:area_law_exponent_gaps}
    For every $p\in\mathbb R$, there exists $B>0$, depending only on $p$ and
    $\delta_0$, such that for every finite set $F$ of nonnegative integers,
    \begin{align}
        \sum_{k\in F}(k+1)^p2^{-\delta_0 k/2} & \le B.
    \end{align}
\end{lemma}
\begin{proof}\leanok
    \uses{thm:area_law_polynomial_summability}
    Let $S\ge0$ be the sum of the series $\sum_{k=0}^\infty w_p(k)$ and set
    $B=S+1>0$. Since $w_p(k)\ge0$, every finite partial sum is at most
    $S<B$.
\end{proof}

\section{Primary birth regions and adjacent fine meshes}
\label{sec:area_law_primary_regions}

Retain the half-open cells $C_{k,z}$, the layer $D_k$, and its finite
index set $I_k$ from Definition~\ref{def:area_law_dyadic_layers}.
The origin $o\in\mathbb R^2$ is arbitrary, the endpoint set
$Z\subset\mathbb Z^2$ is finite, and the neighborhood radius $C_0$ is
a nonnegative integer. All distances and diameters use the sup metric
on $\mathbb R^2$.

The primary regions in \cite[Section~11]{OpenAI2026AreaLaw} are obtained
by intersecting a pitch interior with an actual layer and taking its
closure. All fragments in one pitch interior receive the same primary
identifier, even when their union is disconnected.

% Source: 10-geometry.tex, lines 212--218 and 237--249,
% geometry:primary-pieces. These entries describe individual primary
% birth regions and their fragments; they do not state the full partition
% or the scanner-template conditions of prop:two-families.

\begin{definition}[Pitch interiors and primary birth regions]
    \label{def:area_law_primary_regions}
    \lean{TNLean.PEPS.AreaLaw.Geometry.pitchInterior,
        TNLean.PEPS.AreaLaw.Geometry.primaryFragment,
        TNLean.PEPS.AreaLaw.Geometry.primaryFragmentIndices,
        TNLean.PEPS.AreaLaw.Geometry.primaryBirthRegion}
    \leanok
    \uses{def:area_law_dyadic_cells,def:area_law_dyadic_layers}
    Fix nonnegative integer indices $k,\ell,p$, put
    $r=2^k$, $t=2^\ell$, and $s=2^p$, and choose integer shifts $a,b$.
    For every $j\in\mathbb Z^2$, set
    \begin{align}
        v_j     & =o+t(a,b)+sj,                                \\
        U_j     & =(v_{j,1}+t,v_{j,1}+s)
        \times(v_{j,2}+t,v_{j,2}+s),                           \\
        F_{j,z} & =\overline{C_{k,z}\cap U_j},                 \\
        J_j     & =\{z\in I_k:C_{k,z}\cap U_j\ne\varnothing\}, \\
        R_j     & =\overline{D_k\cap U_j}.
    \end{align}
    Here $U_j$ is the open pitch interior, $F_{j,z}$ is a closed cell
    fragment, and $R_j$ is the corresponding primary birth region.
    These definitions also allow empty interiors and empty fragments.
    In the geometric construction, $\ell=\ell_k$, $p=p_k$, and
    $a,b$ represent the selected residue classes modulo $s/t$.
    In the picture below, with $s=4t$ and $r=2t$, the open pitch interiors
    $U_j$ and $U_{j'}$, for $j'=j+(1,0)$, share one vertical interval and
    are separated by the vertical strip of width $t$ between them, and each
    scale-$r$ cell of the layer
    contributes one closed fragment to the primary region of each pitch
    interior it meets.
    % Figure area_law_primary_regions.
% Owner: chapter/ch24_peps_area_law_primary_regions.tex (def:area_law_primary_regions).
% Used by: gapped-2d-notes AL App. A 2 (primaries and belts), Fig. 8(a).
% Source: 10-geometry.tex, lines 212--249 and the figure
% geometry:layer-picture, lines 253--291 (D_k, pitch s_k, strips t_k).
% Specialization drawn: s = 4t, r = 2t, shifts a = b = 0; rows run with
% the second coordinate increasing upward, columns with the first.
% Ink-to-index: every dot is one closed fine square of side t. The grey
% tinted comb is the belt: the union of the strips [v_{j,i}, v_{j,i}+t]
% removed from each pitch square (lower and left sides). The red enclosures
% are the open pitch interiors U_j and U_{j'}. Each dashed square is one
% scale-r cell C_{k,z} of D_k that straddles a strip; the blue tinted sets
% are the fragments F_{j,z} = closure(C_{k,z} cap U_j). The label R_j is
% inscribed in the whole-cell fragment and R_{j'} hangs on its fragment's
% hull, both in the layer hue. The two blue fragments inside U_j are
% disconnected but form the single region R_j (all fragments in one pitch
% interior make one primary), so R_j is labelled once: every station beside
% the lower fragment touches the dashed straddling cell or a dot. The
% upper-left fragment of R_j is a whole cell whose dashed outline coincides
% with it and is omitted.
% Colours follow the notes' Section 4.1 table (hue-only redeclarations):
% belt grey (context), layer blue (as D_k in area_law_dyadic_layers).
% Not to scale: the source has t << r << s. No strings, legs or
% contractions; empty boundary signature.
\begin{center}
    \begin{tenkz}[size=s, rows={ket,ket,ket,ket}, cols=8, bonds=none]
        \tndeclare{species}{belt}{hue=source:gray}
        \tndeclare{species}{layer}{hue=source:blue!65!black}
        \tnmark[form=enclosure, species=belt, tint]{(1,1) .. (4,8) - (1,2) .. (3,4) - (1,6) .. (3,8)}{}
        \tnmark[form=enclosure, species=layer, tint]{(1,3) .. (2,4)}{}
        \tnmark[form=enclosure, species=layer, tint]{(3,3) .. (3,4)}{}
        \tnmark[form=enclosure, species=layer, tint, label pos=0]{(1,6) .. (2,6)}{$R_{j'}$}
        \tnmark[form=enclosure, species=complement]{(3,3) .. (4,4)}{}
        \tnmark[form=enclosure, species=complement]{(1,5) .. (2,6)}{}
        \tnmark[form=enclosure, species=selected, label pos=n]{(1,2) .. (3,4)}{$U_j$}
        \tnmark[form=enclosure, species=selected, label pos=n]{(1,6) .. (3,8)}{$U_{j'}$}
        \tnmark[form=label, species=layer, label pos=e]{midway (1,3) and (2,3)}{$R_j$}
    \end{tenkz}
\end{center}

\end{definition}

\begin{theorem}[Finite rectangular decomposition]
    \label{thm:area_law_primary_fragment_union}
    \lean{TNLean.PEPS.AreaLaw.Geometry.primaryBirthRegion_eq_iUnion_primaryFragment}
    \leanok
    \uses{def:area_law_primary_regions}
    For every choice of the preceding data,
    \begin{align}
        R_j & =\bigcup_{z\in J_j}F_{j,z}.
    \end{align}
\end{theorem}
\begin{proof}\leanok
    \uses{thm:area_law_dyadic_nesting}
    The exact cell union gives $D_k=\bigcup_{z\in I_k}C_{k,z}$.
    Intersect this finite union with $U_j$ and commute closure with
    the finite union. Empty intersections have empty closures, so
    retaining precisely the indices in $J_j$ leaves the union unchanged.
\end{proof}

\begin{theorem}[Closed cell fragments]
    \label{thm:area_law_primary_fragment_rectangles}
    \lean{TNLean.PEPS.AreaLaw.Geometry.primaryFragment_eq_closedRectangle,
        TNLean.PEPS.AreaLaw.Geometry.primaryFragment_subset_closure_dyadicCell,
        TNLean.PEPS.AreaLaw.Geometry.primaryFragment_dist_le,
        TNLean.PEPS.AreaLaw.Geometry.primaryFragment_diam_le}
    \leanok
    \uses{def:area_law_primary_regions}
    For $i\in\{1,2\}$, write
    $\alpha_{i,z}=o_i+rz_i$ and $\beta_{i,z}=o_i+r(z_i+1)$.
    If $C_{k,z}\cap U_j\ne\varnothing$, then
    \begin{align}
        F_{j,z}
         & =\prod_{i=1}^2
        [\max\{\alpha_{i,z},v_{j,i}+t\},
        \min\{\beta_{i,z},v_{j,i}+s\}].
    \end{align}
    Both coordinate intervals have strictly positive length. Every
    fragment, including an empty one, satisfies the following containment
    and diameter bounds. The distance bound holds for arbitrary points
    $x,y\in F_{j,z}$.
    \begin{align}
        F_{j,z}                            & \subseteq\overline{C_{k,z}}, \\
        \|x-y\|_\infty                     & \le r,                       \\
        \operatorname{diam}_\infty F_{j,z} & \le r.
    \end{align}
\end{theorem}
\begin{proof}\leanok
    \uses{thm:area_law_dyadic_cell_distance}
    A point of the nonempty intersection lies below both upper endpoints
    in each coordinate. It lies at or above the cell's lower endpoint
    and strictly above the pitch interior's lower endpoint. Hence the
    maximum lower endpoint is strictly below the minimum upper endpoint.
    The coordinate intersection contains the open interval between these
    endpoints and is contained in their closed interval. Taking closures
    therefore gives that closed interval. Closure commutes with products,
    proving the rectangle formula.

    Containment in the closed cell follows from monotonicity of closure.
    Apply the closed-cell distance bound with identical cell indices and
    index radius zero. This gives the pointwise distance bound, and hence
    the diameter bound.
\end{proof}

\begin{theorem}[Primary diameter and separation]
    \label{thm:area_law_primary_diameter_separation}
    \lean{TNLean.PEPS.AreaLaw.Geometry.primaryBirthRegion_subset_closure_pitchInterior,
        TNLean.PEPS.AreaLaw.Geometry.primaryBirthRegion_dist_le,
        TNLean.PEPS.AreaLaw.Geometry.primaryBirthRegion_diam_le,
        TNLean.PEPS.AreaLaw.Geometry.primaryBirthRegion_dist_separation}
    \leanok
    \uses{def:area_law_primary_regions}
    For every primary index $j$, the following containment and diameter
    bounds hold. The distance bound holds for arbitrary points $x,y\in R_j$.
    \begin{align}
        R_j                            & \subseteq\overline{U_j}, \\
        \|x-y\|_\infty                 & \le s,                   \\
        \operatorname{diam}_\infty R_j & \le s.
    \end{align}
    For distinct indices $i,j$, every $x\in R_i$ and $y\in R_j$ satisfy
    \begin{align}
        \|x-y\|_\infty & \ge t.
    \end{align}
\end{theorem}
\begin{proof}\leanok
    \uses{def:area_law_primary_regions}
    Monotonicity of closure gives $R_j\subseteq\overline{U_j}$.
    This closure is contained in the product of the closed intervals
    $[v_{j,q}+t,v_{j,q}+s]$, for $q=1,2$. Two points in this product
    have coordinate differences at most $s$, proving the distance and
    diameter bounds.

    Distinct integer indices differ in at least one coordinate. Suppose
    $i_q<j_q$. The upper endpoint for index $i_q$ is at most
    $o_q+t(a,b)_q+sj_q$, while the lower endpoint for index $j_q$ is
    $o_q+t(a,b)_q+sj_q+t$. Thus the two coordinate intervals are
    separated by at least $t$. Taking the maximum coordinate distance
    proves the separation claim; the other ordering is symmetric.
\end{proof}

\begin{theorem}[Number of fragments of one primary]
    \label{thm:area_law_primary_fragment_count}
    \lean{TNLean.PEPS.AreaLaw.Geometry.card_primaryFragmentIndices_le}
    \leanok
    \uses{def:area_law_primary_regions}
    If $k\le p$, then every primary index satisfies
    \begin{align}
        |J_j| & \le\bigl(2^{p-k}+2\bigr)^2
        =\bigl(s/r+2\bigr)^2.
    \end{align}
    This bound is independent of the endpoint set, neighborhood radius,
    origin, and selected shifts.
\end{theorem}
\begin{proof}\leanok
    \uses{def:area_law_primary_regions}
    Put $N=2^{p-k}$, so $s=Nr$, and let
    $q_i=\lfloor(v_{j,i}-o_i)/r\rfloor$. If a cell indexed by $z$ meets
    $U_j$, a point in that intersection lies above $v_{j,i}$ and below
    $v_{j,i}+s$. The half-open cell inequalities and the defining floor
    inequalities consequently give, for each coordinate $i\in\{1,2\}$,
    \begin{align}
        q_i & \le z_i\le q_i+N.
    \end{align}
    Thus $J_j$ lies in an integer rectangle with $N+1$ indices in each
    coordinate. Its cardinality is at most $(N+1)^2$, which is at most
    $(N+2)^2$.
\end{proof}

% Source: 10-geometry.tex, lines 299--306 and 352--356,
% the adjacent-mesh arithmetic used in geometry:initial-stars.
\begin{theorem}[Adjacent fine meshes]
    \label{thm:area_law_adjacent_fine_meshes}
    \lean{TNLean.PEPS.AreaLaw.Geometry.fineScaleIndex_mono,
        TNLean.PEPS.AreaLaw.Geometry.fineScaleIndex_succ,
        TNLean.PEPS.AreaLaw.Geometry.fineScale_side_succ}
    \leanok
    \uses{def:area_law_fine_belt_scales}
    The fixed fine exponents $\ell_k=\lfloor\zeta k\rfloor$ are
    nondecreasing. For every nonnegative integer $k$,
    \begin{align}
        \ell_{k+1} & \in\{\ell_k,\ell_k+1\}, \\
        t_{k+1}    & \in\{t_k,2t_k\}.
    \end{align}
\end{theorem}
\begin{proof}\leanok
    \uses{thm:area_law_exponent_gaps}
    Since $0\le\zeta\le1$, multiplication by $\zeta$ and taking floors
    preserve the order of the layer indices. Also
    $\zeta(k+1)\le\zeta k+1$, and adding one commutes with taking the
    floor. Therefore $\ell_k\le\ell_{k+1}\le\ell_k+1$.
    Exponentiating the two possible integer values gives the side-length
    assertion.
\end{proof}

\section{Birth templates and scanner scales of primaries}
\label{sec:area_law_primary_birth_templates}

Retain the primary birth regions $R_j$, fragments $F_{j,z}$, and index sets
$J_j$ (Definition~\ref{def:area_law_primary_regions}) at layer $k$ with
$r=2^k$, fine side $t_k$, and belt pitch $s_k$
(Definition~\ref{def:area_law_fine_belt_scales}). A primary is present iff
$J_j\ne\varnothing$, in which case \cite[Section~11]{OpenAI2026AreaLaw} assigns
its scanner scale and birth template as follows.

% Source: 10-geometry.tex, lines 709--729 (geometry:primary-scale and
% geometry:primary-limits) and lines 756--774 (order of the constants).

\begin{definition}[Primary scanner scale and birth template]
    \label{def:area_law_primary_birth_template}
    \lean{TNLean.PEPS.AreaLaw.Geometry.primaryScannerScale,
        TNLean.PEPS.AreaLaw.Geometry.primaryBirthTemplate}
    \leanok
    \uses{def:area_law_primary_regions,def:area_law_fine_belt_scales,
        def:area_law_template}
    For a constant $C_P$, define the primary scanner scale at layer $k$ by
    \begin{align}
        n_k & =\bigl\lceil C_P r^{1+2\delta_0}\bigr\rceil.
        \label{eq:area_law_primary_scanner_scale}
    \end{align}
    For $J_j\ne\varnothing$ and $n\in\mathbb N_{\ge1}$ with
    $C_{\rm tpl}|J_j|(r+1)\le n$ and $2^p\le n$, the primary birth template is
    the template $T$ (Definition~\ref{def:area_law_template}) with parameters
    $(n,s_0)=(n,r)$, pieces $\{F_{j,z}\}_{z\in J_j}$, and points
    $\bigcup_{z\in J_j}(F_{j,z}\cap\mathbb Z^2)$.
\end{definition}

\begin{lemma}[Points and pieces of a primary birth template]
    \label{thm:area_law_primary_birth_template_points}
    \lean{TNLean.PEPS.AreaLaw.Geometry.primaryBirthTemplate_points,
        TNLean.PEPS.AreaLaw.Geometry.primaryBirthTemplate_pieceCount,
        TNLean.PEPS.AreaLaw.Geometry.card_primaryFragmentIndices_le_rpow,
        TNLean.PEPS.AreaLaw.Geometry.card_primaryBirthTemplate_pieces_le}
    \leanok
    The primary birth template $T$ of Definition~\ref{def:area_law_primary_birth_template}
    has $N_0=|J_j|$ pieces with $\operatorname{diam}_\infty(F_{j,z})\le r$,
    $\operatorname{diam}_\infty(T)\le 2^p$, and (at $p=p_k$)
    \begin{align}
        T   & =R_j\cap\mathbb Z^2, &
        N_0 & \le 9r^{2\delta_0}.
        \label{eq:area_law_primary_birth_points_count}
    \end{align}
\end{lemma}
\begin{proof}\leanok
    \uses{def:area_law_primary_birth_template,
        thm:area_law_primary_fragment_union,
        thm:area_law_primary_fragment_rectangles,
        thm:area_law_primary_diameter_separation,
        thm:area_law_primary_fragment_count,
        thm:area_law_fine_belt_scales}
    Each $F_{j,z}$ ($z\in J_j$) has endpoints of the form $o_i+m$ ($m\in\mathbb Z$)
    with strict inequality, so $F_{j,z}\cap\mathbb Z^2\ne\varnothing$; since
    $R_j=\bigcup_{z\in J_j}F_{j,z}$ (Lemma~\ref{thm:area_law_primary_fragment_union}),
    $T=R_j\cap\mathbb Z^2$, and Lemmas~\ref{thm:area_law_primary_fragment_rectangles}
    and~\ref{thm:area_law_primary_diameter_separation} give the diameter bounds.
    With $m=p_k-k\ge0$, $p_k\le(1+\delta_0)k$
    (Lemma~\ref{thm:area_law_fine_belt_scales}) gives $1\le 2^m\le r^{\delta_0}$,
    so Lemma~\ref{thm:area_law_primary_fragment_count} yields
    \begin{align}
        N_0 & \le(2^m+2)^2\le9\cdot4^m\le9r^{2\delta_0}.
        \notag
    \end{align}
\end{proof}

\begin{lemma}[Choice of the primary scale constant]
    \label{thm:area_law_primary_scanner_constant}
    \lean{TNLean.PEPS.AreaLaw.Geometry.primaryBirthTemplate_scale,
        TNLean.PEPS.AreaLaw.Geometry.two_pow_pitchScaleIndex_le_primaryScannerScale,
        TNLean.PEPS.AreaLaw.Geometry.exists_primaryScannerConstant}
    \leanok
    For any $C_{\rm tpl}$, setting $C_P=\max\{18C_{\rm tpl},1\}$ in~(\ref{eq:area_law_primary_scanner_scale})
    ensures that every primary at belt pitch $p_k$ satisfies
    \begin{align}
        n_k                          & \ge C_{\rm tpl}N_0(r+1), &
        \operatorname{diam}_\infty T & \le s_k\le n_k,
        \label{eq:area_law_primary_scanner_constant}
    \end{align}
    so $T$ is a template of Definition~\ref{def:area_law_template} with
    parameters $(n_k,r)$.
\end{lemma}
\begin{proof}\leanok
    \uses{thm:area_law_primary_birth_template_points,
        def:area_law_primary_birth_template}
    For $C_{\rm tpl}>0$ and $r\ge1$,~(\ref{eq:area_law_primary_birth_points_count})
    gives
    \begin{align}
        C_{\rm tpl}N_0(r+1)
         & \le C_{\rm tpl}\cdot9r^{2\delta_0}\cdot2r
        =18C_{\rm tpl}r^{1+2\delta_0}\le n_k,
        \notag
    \end{align}
    and $C_P\ge1$ gives $s_k=2^{p_k}\le r^{1+\delta_0}\le r^{1+2\delta_0}\le n_k$.
\end{proof}

\begin{lemma}[Cut clearance of a primary birth template]
    \label{thm:area_law_primary_birth_template_clearance}
    \lean{TNLean.PEPS.AreaLaw.Geometry.primaryBirthTemplate_isSeparated}
    \leanok
    If $4D_0<C_0$ (for example $C_0>100(D_0+1)$ as in
    \cite[Section~11]{OpenAI2026AreaLaw}), then every primary birth template $T$
    satisfies
    \begin{align}
        \operatorname{dist}_\infty(T,Z) & >4D_0s_0.
        \label{eq:area_law_primary_birth_clearance}
    \end{align}
\end{lemma}
\begin{proof}\leanok
    \uses{thm:area_law_primary_birth_template_points,
        thm:area_law_dyadic_layer_distance}
    Since $T\subseteq R_j\subseteq\overline{D_k}$
    (Lemma~\ref{thm:area_law_primary_birth_template_points}),
    Lemma~\ref{thm:area_law_dyadic_layer_distance} gives
    $\operatorname{dist}_\infty(T,Z)\ge C_0r>4D_0r=4D_0s_0$.
\end{proof}

\begin{lemma}[Primary scale limits]
    \label{thm:area_law_primary_scale_limits}
    \lean{TNLean.PEPS.AreaLaw.Geometry.eventually_primaryScaleLimits,
        TNLean.PEPS.AreaLaw.Geometry.le_primaryScannerScale}
    \leanok
    For any $C_P$ and $\varepsilon>0$, all sufficiently large $k$ (uniformly
    over primaries) satisfy
    \begin{align}
        \bigl\lfloor n_k^{1-\ell}\bigr\rfloor & \le\varepsilon r,   &
        n_k^\beta                             & \le\varepsilon t_k, &
        n_k^\beta                             & \le\varepsilon r,
        \label{eq:area_law_primary_scale_limits}
    \end{align}
    and if $C_P>0$, then for every $n_*$, $n_k\ge n_*$ for all sufficiently
    large $k$.
\end{lemma}
\begin{proof}\leanok
    \uses{def:area_law_primary_birth_template,thm:area_law_exponent_gaps}
    With $e=1+2\delta_0$ and $C=\max\{C_P,0\}+1$,~(\ref{eq:area_law_primary_scanner_scale})
    gives $n_k\le Cr^e$. Since $e(1-\ell)<1$, $e\beta<\zeta$
    (Lemma~\ref{thm:area_law_exponent_gaps}), and $r^\zeta/2\le t_k\le r$,
    $Cr^{-(b-a)}\to0$ ($a<b$) yields~(\ref{eq:area_law_primary_scale_limits});
    for $C_P>0$, $n_k\ge C_P r\to\infty$.
\end{proof}

\section{The number of primary identifiers in a layer}
\label{sec:area_law_primary_count}

Retain the origin, layer and pitch interiors from
Definition~\ref{def:area_law_primary_regions}. Write $I_k$ for the
layer-cell index set, $r=2^k$, $t=2^\ell$, $s=2^p$, and
$o'=o+t(a,b)$. The count below is the primary count in
\cite[Section~11]{OpenAI2026AreaLaw}.

% Source: 10-geometry.tex, lines 212--218 and 237--238.
% The entries count the actual primary interiors in a single layer;
% they do not assert the remaining two-family or repair properties.

\begin{definition}[Primary pitch indices]
    \label{def:area_law_primary_pitch_indices}
    \lean{TNLean.PEPS.AreaLaw.Geometry.primaryPitchIndices}
    \leanok
    \uses{def:area_law_primary_regions,def:area_law_dyadic_cells}
    For every $z\in I_k$, let $L_z$ and $H_z$ be the lower and upper
    corners of the closed cell of side $r$. For each coordinate $i$,
    define
    \begin{align}
        q_{z,i}^- & =\left\lfloor\frac{L_{z,i}-o'_i}{s}\right\rfloor, \\
        q_{z,i}^+ & =\left\lfloor\frac{H_{z,i}-o'_i}{s}\right\rfloor.
    \end{align}
    Let $B_z$ be the integer rectangle
    $[q_{z,1}^-,q_{z,1}^+]\times[q_{z,2}^-,q_{z,2}^+]$.
    The finite set of primary pitch indices is
    \begin{align}
        \mathcal P_k
         & =\left\{j\in\bigcup_{z\in I_k}B_z:
        D_k\cap U_j\ne\varnothing\right\}.
    \end{align}
\end{definition}

\begin{theorem}[Exact primary membership]
    \label{thm:area_law_primary_pitch_membership}
    \lean{TNLean.PEPS.AreaLaw.Geometry.mem_primaryPitchIndices_iff}
    \leanok
    \uses{def:area_law_primary_pitch_indices}
    For every choice of the nonnegative scale indices and every
    $j\in\mathbb Z^2$,
    \begin{align}
        j\in\mathcal P_k & \quad\Longleftrightarrow\quad
        D_k\cap U_j\ne\varnothing.
    \end{align}
\end{theorem}
\begin{proof}\leanok
    \uses{thm:area_law_dyadic_nesting}
    Only the reverse implication needs proof. A point $x\in D_k\cap U_j$
    belongs to a cell $C_{k,z}$ with $z\in I_k$. Since $L_{z,i}\le x_i\le H_{z,i}$,
    monotonicity of the floor gives
    $q_{z,i}^-\le\lfloor(x_i-o'_i)/s\rfloor\le q_{z,i}^+$.
    From $x\in U_j$ we have $j_i<(x_i-o'_i)/s<j_i+1$, so
    $\lfloor(x_i-o'_i)/s\rfloor=j_i$ and $j\in B_z$.
\end{proof}

\begin{theorem}[Primary count from the layer-cell count]
    \label{thm:area_law_primary_pitch_count}
    \lean{TNLean.PEPS.AreaLaw.Geometry.card_primaryPitchIndices_le}
    \leanok
    \uses{def:area_law_primary_pitch_indices}
    If $k\le p$, then
    \begin{align}
        |\mathcal P_k| & \le4|I_k|.
    \end{align}
\end{theorem}
\begin{proof}\leanok
    \uses{def:area_law_dyadic_cells}
    Since $k\le p$, we have $r\le s$. In each coordinate,
    $H_{z,i}=L_{z,i}+r\le L_{z,i}+s$, so
    \begin{align}
        q_{z,i}^- & \le q_{z,i}^+\le q_{z,i}^-+1.
    \end{align}
    Hence $|B_z|\le4$ for every $z\in I_k$, and summing over $z\in I_k$ gives
    $|\mathcal P_k|\le4|I_k|$.
\end{proof}

\begin{theorem}[Primary count from the cut boundary]
    \label{thm:area_law_primary_boundary_count}
    \lean{TNLean.PEPS.AreaLaw.Geometry.card_primaryPitchIndices_boundary_le}
    \leanok
    \uses{def:area_law_primary_pitch_indices,def:area_law_hamiltonian}
    Take $Z$ to be the endpoint set of the crossing edges of a cut $A$
    in a finite domain $\Lambda$. If $k\le p$, then
    \begin{align}
        |\mathcal P_k| & \le32(2C_0+1)^2|\partial_\Lambda A|.
    \end{align}
\end{theorem}
\begin{proof}\leanok
    \uses{thm:area_law_primary_pitch_count,thm:area_law_dyadic_layer_count}
    Combine $|\mathcal P_k|\le4|I_k|$ with
    $|I_k|\le8(2C_0+1)^2|\partial_\Lambda A|$.
\end{proof}

\section{The marked points of belt cells}
\label{sec:area_law_belt_marks}

The belt construction marks each cell at its center, corners and side
midpoints. We first count these points for a finite family of cells
at one scale, as in \cite[Section~11]{OpenAI2026AreaLaw}.
Coincident points contribute only once to the union.

% Source: 10-geometry.tex, lines 325--330; the final estimate also uses
% geometry:belt-count, lines 220--228. Minimum incident scales and the
% initial-star statement are separate results.

\begin{definition}[The nine marks of one cell]
    \label{def:area_law_belt_cell_marks}
    \lean{TNLean.PEPS.AreaLaw.Geometry.beltCellMarks}
    \leanok
    \uses{def:area_law_dyadic_cells}
    For a cell $C_{\ell,z}=o+t(z+[0,1)^2)$, where $t=2^\ell$, put
    \begin{align}
        M_{\ell,z}
         & =\left\{o+t\left(z+\frac e2\right):
        e\in\{0,1,2\}^2\right\}.
    \end{align}
    These are its four corners, four side midpoints and center.
\end{definition}

\begin{definition}[The union of cell marks]
    \label{def:area_law_belt_marks}
    \lean{TNLean.PEPS.AreaLaw.Geometry.beltMarks}
    \leanok
    \uses{def:area_law_belt_cell_marks}
    For any finite cell-index set $F\subset\mathbb Z^2$, define
    \begin{align}
        M_\ell(F) & =\bigcup_{z\in F}M_{\ell,z}.
    \end{align}
\end{definition}

\begin{theorem}[Marks lie in their closed cell]
    \label{thm:area_law_belt_marks_closed_cell}
    \lean{TNLean.PEPS.AreaLaw.Geometry.beltCellMarks_subset_closure_dyadicCell}
    \leanok
    \uses{def:area_law_belt_cell_marks}
    Every point of $M_{\ell,z}$ belongs to $\overline{C_{\ell,z}}$.
\end{theorem}
\begin{proof}\leanok
    \uses{thm:area_law_dyadic_closures}
    Each coordinate of $e/2$ lies in $[0,1]$, so both coordinates
    of each mark lie in the corresponding closed cell intervals.
\end{proof}

\begin{theorem}[Exact cell-mark count]
    \label{thm:area_law_belt_cell_mark_count}
    \lean{TNLean.PEPS.AreaLaw.Geometry.card_beltCellMarks}
    \leanok
    \uses{def:area_law_belt_cell_marks}
    For every origin $o$, scale $\ell$ and cell index $z$, $|M_{\ell,z}|=9$.
\end{theorem}
\begin{proof}\leanok
    Since $2^\ell>0$, distinct elements of $\{0,1,2\}^2$ yield distinct points,
    and there are nine such elements.
\end{proof}

\begin{theorem}[Count of the union of marks]
    \label{thm:area_law_belt_mark_count}
    \lean{TNLean.PEPS.AreaLaw.Geometry.card_beltMarks_le}
    \leanok
    \uses{def:area_law_belt_marks}
    For every finite cell-index set $F\subset\mathbb Z^2$,
    \begin{align}
        |M_\ell(F)| & \le9|F|.
    \end{align}
\end{theorem}
\begin{proof}\leanok
    \uses{thm:area_law_belt_cell_mark_count}
    The cardinality of a finite union is at most the sum of the
    cardinalities of its members, each of which has nine points by
    Theorem~\ref{thm:area_law_belt_cell_mark_count}.
\end{proof}

\begin{theorem}[Sparse shifts bound the belt marks]
    \label{thm:area_law_sparse_belt_marks}
    \lean{TNLean.PEPS.AreaLaw.Geometry.exists_sparse_dyadic_belt_marks_shift}
    \leanok
    \uses{def:area_law_belt_marks,def:area_law_belt_cells,
        def:area_law_fine_layer_indices,def:area_law_fine_belt_scales}
    Let $F_k(a,b)$ be the belt-cell indices in the refined
    layer for coordinate residues $a,b$. There is a pair
    of residues $(a,b)$ such that
    \begin{align}
        |M_{\ell_k}(F_k(a,b))|
         & \le576(2C_0+1)^2|\partial_\Lambda A|\,2^{-\delta_0 k/2}.
    \end{align}
    The estimate is uniform in the finite domain, cut, origin and scale.
\end{theorem}
\begin{proof}\leanok
    \uses{thm:area_law_belt_mark_count,thm:area_law_sparse_dyadic_belts}
    Choose the pair of residues given by Theorem~\ref{thm:area_law_sparse_dyadic_belts}
    and apply Theorem~\ref{thm:area_law_belt_mark_count}. The coefficient is
    $9\cdot64=576$.
\end{proof}

\section{Separation on the translated quarter mesh}
\label{sec:area_law_quarter_mesh}

The numerical estimates used in the initial-star argument come from a
common translated lattice. We record them before constructing
the active boundaries; see \cite[Section~11]{OpenAI2026AreaLaw}.

% Source: 10-geometry.tex, geometry:initial-stars, lines 352--359.
% These estimates concern actual marked points and allowed affine lines.
% Local layer exclusion is included below; active fan interfaces and
% isolated stars remain separate.

\begin{theorem}[Nonadjacent layers avoid fine neighborhoods]
    \label{thm:area_law_local_layer_exclusion}
    \lean{TNLean.PEPS.AreaLaw.Geometry.dyadicLayer_dist_nonadjacent_fineScale}
    \leanok
    \uses{def:area_law_dyadic_layers,def:area_law_fine_belt_scales}
    Let $C_0\ge2$ and $k\ge50{,}000{,}000$, and put $t_k=2^{\ell_k}$.
    If $x\in\overline{D_k}$, $y\in\overline{D_h}$, and the layer
    indices differ by at least two, then
    \begin{align}
        \|x-y\|_\infty & \ge16t_k.
    \end{align}
    The threshold is uniform in the finite endpoint set and origin.
\end{theorem}
\begin{proof}\leanok
    \uses{thm:area_law_dyadic_nonadjacent,thm:area_law_scale_side_ratios}
    The scale-ratio estimate with exponent gap five gives $32t_k\le2^k$.
    The nonadjacent-layer bound gives
    $\|x-y\|_\infty\ge(C_0-1)2^{\max(k,h)}/2\ge2^k/2\ge16t_k$.
\end{proof}

\begin{theorem}[Only neighboring layers occur nearby]
    \label{thm:area_law_local_layer_indices}
    \lean{TNLean.PEPS.AreaLaw.Geometry.dyadicLayer_nearby_indices}
    \leanok
    \uses{def:area_law_dyadic_layers,def:area_law_fine_belt_scales}
    Under the same conditions on $C_0$ and $k$, if
    $x\in\overline{D_k}$, $y\in\overline{D_h}$ and
    $\|x-y\|_\infty<10t_k$, then
    \begin{align}
        h & \le k+1, & k & \le h+1.
    \end{align}
\end{theorem}
\begin{proof}\leanok
    \uses{thm:area_law_local_layer_exclusion}
    A difference of at least two would give distance at least $16t_k$,
    contradicting the strict upper bound since $t_k>0$.
\end{proof}

\begin{definition}[Translated planar mesh]
    \label{def:area_law_affine_mesh}
    \lean{TNLean.PEPS.AreaLaw.Geometry.affineMesh}
    \leanok
    For $o\in\mathbb R^2$ and $q\in\mathbb R$, define
    \begin{align}
        \mathcal M(o,q) & =\{o+qz:z\in\mathbb Z^2\}.
    \end{align}
\end{definition}

\begin{theorem}[Belt marks lie on the quarter mesh]
    \label{thm:area_law_belt_marks_quarter_mesh}
    \lean{TNLean.PEPS.AreaLaw.Geometry.beltMarks_subset_affineMesh}
    \leanok
    \uses{def:area_law_affine_mesh,def:area_law_belt_marks}
    Let $j,\ell$ be nonnegative integers with $\ell\le j+1$.
    For every finite cell-index set $F$,
    \begin{align}
        M_j(F) & \subseteq\mathcal M(o,2^\ell/4).
    \end{align}
    In particular, cells of sides $t/2$, $t$ and $2t$, where
    $t=2^\ell$ and $\ell\ge1$, have all their marks on the same
    mesh $\mathcal M(o,t/4)$.
\end{theorem}
\begin{proof}\leanok
    A mark of a cell of side $2^j$ is $o+2^j(z+e/2)$, where
    $e\in\{0,1,2\}^2$. Since $\ell\le j+1$, the vector
    $2^{j+1-\ell}(2z+e)$ is integral, and the mark equals
    $o+(2^\ell/4)2^{j+1-\ell}(2z+e)$.
\end{proof}

\begin{theorem}[Distinct mesh points are separated]
    \label{thm:area_law_affine_mesh_point_separation}
    \lean{TNLean.PEPS.AreaLaw.Geometry.affineMesh_dist_ge}
    \leanok
    \uses{def:area_law_affine_mesh}
    If $q>0$ and $x,y\in\mathcal M(o,q)$ are distinct, then
    \begin{align}
        \|x-y\|_\infty & \ge q.
    \end{align}
\end{theorem}
\begin{proof}\leanok
    At least one coordinate of the integer vectors indexing $x$ and
    $y$ differs. Its difference has absolute value at least one, so
    the corresponding coordinate difference of $x$ and $y$ has
    absolute value at least $q$.
\end{proof}

\begin{theorem}[Clearance from a nonincident allowed line]
    \label{thm:area_law_affine_mesh_line_clearance}
    \lean{TNLean.PEPS.AreaLaw.Geometry.affineMesh_line_dist_ge}
    \leanok
    \uses{def:area_law_affine_mesh,def:area_law_template}
    Suppose $q>0$ and $u,v,x\in\mathcal M(o,q)$. Let $L$ be the
    affine span of $u,v$. Suppose either $u=v$, or $L$ has slope
    $0$, $\infty$, $1$ or $-1$. If $x\notin L$, then
    \begin{align}
        \|x-y\|_\infty & \ge q/2\qquad(y\in L).
    \end{align}
    Thus on a quarter mesh the clearance is at least $t/8$.
\end{theorem}
\begin{proof}\leanok
    \uses{thm:area_law_affine_mesh_point_separation}
    If $u=v$, use the point-separation bound. Otherwise, the supporting
    line has a constant coordinate, coordinate sum or coordinate
    difference. Its value at a missing mesh point differs from the
    line's value by a nonzero integer multiple of $q$. A coordinate
    difference is bounded by the sup distance; a difference of sums
    or differences is bounded by twice that distance. The result follows.
\end{proof}

\begin{theorem}[Clearance from a nonincident allowed segment]
    \label{thm:area_law_affine_mesh_segment_clearance}
    \lean{TNLean.PEPS.AreaLaw.Geometry.affineMesh_segment_dist_ge_of_not_mem}
    \leanok
    \uses{def:area_law_affine_mesh,def:area_law_template}
    Let $q>0$ and $a,b,p\in\mathcal M(o,q)$. Suppose that the segment
    $[a,b]$ is horizontal, vertical or diagonal, with coincident endpoints
    also permitted. If $p\notin[a,b]$, then every $x\in[a,b]$ satisfies
    $\|p-x\|_\infty\ge q/2$.
    This estimate is used in the isolation arguments for the initial stars
    and recursive repair in \cite[Section~11]{OpenAI2026AreaLaw}.
\end{theorem}
\begin{proof}\leanok
    \uses{thm:area_law_affine_mesh_point_separation,thm:area_law_affine_mesh_line_clearance}
    If $a=b$, then $x=a\ne p$, so mesh-point separation gives the
    stronger bound $q$. Otherwise let $L$ be the supporting line of
    $[a,b]$. If $p\notin L$, apply the line-clearance estimate.
    If $p\in L\setminus[a,b]$, one endpoint $e\in\{a,b\}$ lies
    between $p$ and $x$. Since $p\ne e$, mesh-point separation and
    distance additivity give
    \begin{align}
        q & \le\|p-e\|_\infty
        \le\|p-e\|_\infty+\|e-x\|_\infty
        =\|p-x\|_\infty.
    \end{align}
    In both cases the asserted half-spacing bound follows.
\end{proof}

\subsection{Clipping an incident segment}
\label{subsec:area_law_segment_clipping}

The local repair construction in \cite[Section~11]{OpenAI2026AreaLaw}
intersects the segments incident to a mark with a small closed square.
After normalization, the following lemma determines their endpoints in
any real normed space, and in particular for the planar maximum norm.

\begin{lemma}[Clipping a unit segment]
    \label{thm:area_law_unit_segment_ball_clipping}
    \lean{TNLean.PEPS.AreaLaw.Geometry.segment_inter_closedBall_of_dist_eq_one}
    \leanok
    Let $E$ be a real normed space and let $a,b\in E$ satisfy
    $\|b-a\|=1$. If $0\le\eta\le1$, then
    \begin{align}
        [a,b]\cap\overline B(a,\eta)
         & = [a,(1-\eta)a+\eta b].
    \end{align}
\end{lemma}
\begin{proof}\leanok
    Write a point of $[a,b]$ as $(1-t)a+tb$, where $0\le t\le1$.
    Its distance from $a$ is $t\|b-a\|=t$. It therefore belongs to
    the closed ball precisely when $0\le t\le\eta$.
    These points form the segment on the right-hand side.
\end{proof}

\section{The finite family of initial marks}
\label{sec:area_law_initial_mark_family}

The sparse shifts may be chosen in every distance layer. Their cell counts
have a summable geometric bound, so only finitely many layers contribute
belt cells. We can therefore form the finite union of their actual marked
points and assign to each point its smallest incident cell side; see
\cite[Section~11]{OpenAI2026AreaLaw}.

% Source: 10-geometry.tex, geometry:belt-count, lines 220--233;
% initial marks, lines 325--330; geometry:total-repairs, lines 668--692.
% Only the zero-polynomial summation for initial cells and marks is proved
% here. Point separation is the first assertion of geometry:initial-stars,
% lines 325--339, with its proof at lines 352--359. Repairs remain separate.

\begin{theorem}[Sparse shifts with finite support]
    \label{thm:area_law_finite_sparse_belt_family}
    \lean{TNLean.PEPS.AreaLaw.Geometry.exists_finitely_supported_sparse_belt_shifts}
    \leanok
    \uses{def:area_law_belt_cells,def:area_law_fine_layer_indices,
        def:area_law_fine_belt_scales}
    For every finite domain $\Lambda$, cut $A$, origin $o$ and
    nonnegative integer $C_0$, there are coordinate residues $a_k,b_k$
    at every scale. Write $F_k$ for their actual belt-cell index set.
    These choices satisfy
    \begin{align}
        |F_k| & \le64(2C_0+1)^2|\partial_\Lambda A|\,2^{-\delta_0 k/2}
        \qquad(k\ge0),
    \end{align}
    and the set $\{k:F_k\ne\varnothing\}$ is finite.
\end{theorem}
\begin{proof}\leanok
    \uses{thm:area_law_sparse_dyadic_belts,thm:area_law_polynomial_summability}
    Choose a sparse pair of residues at each scale. The integer counts
    $|F_k|$ are dominated by a fixed multiple of a summable geometric
    sequence. Their real-valued series is therefore summable, so its
    terms tend to zero. Eventually $|F_k|<1$, which forces $F_k$ to be
    empty. Only finitely many indices remain.
\end{proof}

\begin{theorem}[Uniform count of all initial marks]
    \label{thm:area_law_uniform_initial_mark_count}
    \lean{TNLean.PEPS.AreaLaw.Geometry.exists_uniform_initial_mark_bound}
    \leanok
    \uses{def:area_law_belt_marks,def:area_law_belt_cells,
        def:area_law_fine_layer_indices,def:area_law_fine_belt_scales}
    There is a positive constant $B$ such that, for every finite domain,
    cut, origin, nonnegative integer $C_0$ and lower index $k_0$, one can
    choose the residues from the preceding theorem and a finite set $M$
    of actual marks with
    \begin{align}
        x\in M
            & \quad\Longleftrightarrow\quad
        \exists k\ge k_0,\ x\in M_{\ell_k}(F_k),   \\
        |M| & \le B(2C_0+1)^2|\partial_\Lambda A|.
    \end{align}
    The chosen residues retain their per-layer cell estimate and finite
    support. The constant $B$ precedes all the geometric data; the
    support cutoff may depend on those data.
\end{theorem}
\begin{proof}\leanok
    \uses{thm:area_law_finite_sparse_belt_family,thm:area_law_belt_mark_count,
        thm:area_law_polynomial_finite_sums}
    Let $B_0>0$ be the uniform finite-sum budget for the zero polynomial
    exponent, and take $B=576B_0$. Form the union of the actual marks
    over the finite support restricted to $k\ge k_0$. Its membership is
    exactly the stated condition. A finite union has at most the sum of
    the individual cardinalities, and each cell contributes at most
    nine marks. Thus its count is at most
    \begin{align}
        9\sum_{k\ge k_0}|F_k|
         & \le576(2C_0+1)^2|\partial_\Lambda A|
        \sum_{k\in H}2^{-\delta_0 k/2}
        \le B(2C_0+1)^2|\partial_\Lambda A|,
    \end{align}
    where $H$ is the finite retained support.
\end{proof}


\begin{theorem}[Separation of fine-layer marks]
    \label{thm:area_law_fine_layer_mark_separation}
    \lean{TNLean.PEPS.AreaLaw.Geometry.fineLayer_marks_dist_ge}
    \leanok
    \uses{def:area_law_belt_marks,def:area_law_fine_layer_indices,
        def:area_law_fine_belt_scales}
    Fix an origin $o$ and a finite endpoint set $Z$, and form all layers
    from these data with $C_0\ge2$. Let $k,h\ge50{,}000{,}000$. If $x$ and $y$ are
    distinct marks of actual fine-layer cells in layers $k$ and $h$,
    respectively, then
    \begin{align}
        \operatorname{dist}_\infty(x,y)\ge\frac14\max(t_k,t_h),
        \qquad t_j=2^{\ell_j}.
    \end{align}
    Cell-boundary marks are included, and the estimate applies to every
    fine-layer cell, hence to every selected belt cell.
\end{theorem}
\begin{proof}\leanok
    \uses{thm:area_law_local_layer_indices,thm:area_law_adjacent_fine_meshes,
        thm:area_law_belt_marks_quarter_mesh,thm:area_law_affine_mesh_point_separation,
        thm:area_law_fine_cell_containment}
    Suppose that the distance is less than $t_k/4$. Each mark is in the
    closure of its actual layer. Layer locality forces $k\le h+1$.
    Monotonicity and the adjacent-scale estimate give
    $\ell_k\le\ell_h+1$, so both marks lie on the translated mesh of
    spacing $t_k/4$. Distinct mesh points cannot be closer than this
    spacing. Reversing $k$ and $h$ proves both lower bounds and hence
    their maximum.
\end{proof}

\begin{theorem}[Smallest incident sides]
    \label{thm:area_law_least_incident_mark_scales}
    \lean{TNLean.PEPS.AreaLaw.Geometry.exists_least_incident_scales}
    \leanok
    \uses{def:area_law_belt_marks,def:area_law_fine_belt_scales}
    Fix an origin $o$, a lower index $k_0$, and a finite set $F_k$ of
    cell indices at each layer. Let
    \begin{align}
        M=\bigcup_{k\ge k_0}M_{\ell_k}(F_k).
    \end{align}
    There is a positive side $s(v)$ for every $v\in M$, realized by a
    cell indexed by $F_k$ in some layer $k\ge k_0$ that marks $v$.
    Every cell indexed by $F_h$ in a layer $h\ge k_0$ that marks $v$
    has side at least $s(v)$. Neither $M$ nor the set of nonempty layers
    is required to be finite.
\end{theorem}
\begin{proof}\leanok
    \uses{thm:area_law_adjacent_fine_meshes}
    For each $v\in M$, take the least layer index $k(v)\ge k_0$
    with an incident cell, and put $s(v)=2^{\ell_{k(v)}}$.
    The exponent $\ell_k$ is nondecreasing in $k$, so $s(v)$ is a
    smallest incident side. It is realized at the chosen layer and
    is positive.
\end{proof}

\begin{theorem}[Counted and separated initial marks]
    \label{thm:area_law_uniform_separated_initial_marks}
    \lean{TNLean.PEPS.AreaLaw.Geometry.exists_uniform_separated_initial_marks}
    \leanok
    \uses{thm:area_law_uniform_initial_mark_count,def:area_law_belt_marks}
    There is $B>0$ such that, for every finite domain, cut, origin,
    integer $C_0\ge2$ and lower index $k_0\ge50{,}000{,}000$, there are
    sparse residues with the per-layer estimates and finite support,
    the exact finite union $M$ of their actual marks above $k_0$, and a
    positive side $s(v)$ for each $v\in M$, satisfying
    \begin{align}
        |M| & \le B(2C_0+1)^2|\partial_\Lambda A|,     \\
        \operatorname{dist}_\infty(v,w)
            & \ge\frac14\max(s(v),s(w))\qquad(v\ne w).
    \end{align}
    For each $v$, some selected belt cell in a layer $k\ge k_0$ has $v$
    as an actual mark and side $s(v)=t_k$. Every selected belt cell in a layer $h\ge k_0$ that has $v$ as one
    of its actual marks has side at least $s(v)$.
\end{theorem}
\begin{proof}\leanok
    \uses{thm:area_law_uniform_initial_mark_count,
        thm:area_law_least_incident_mark_scales,
        thm:area_law_fine_layer_mark_separation}
    Choose the counted finite union and residues from the preceding
    construction. For each point of $M$, take the least layer index
    admitting an incident selected cell. The fine-cell exponent is
    nondecreasing, so this gives a positive smallest incident side.
    Apply fine-layer mark separation to the two witnessing cells.
    When $M$ is empty the assignment and pairwise condition are empty.
\end{proof}

\section{Layers and meshes near a cell mark}
\label{sec:area_law_near_mark_geometry}

Fix an arbitrary origin $o$, a finite endpoint set $Z\subset\mathbb Z^2$,
and width $C\ge 2$. Let $\ell_k$ be the fine exponent and
$M_{k,z}$ the nine marks of the fine cell indexed by $z$ at
layer $k$. The estimates below establish the local mesh and scale bounds in
\cite[Section~11]{OpenAI2026AreaLaw}.

% Source: 10-geometry.tex, geometry:initial-stars, lines 352--359;
% geometry:layer-distance, lines 179--191; scale comparison, lines 200--207;
% revision adc7f1241b42e322a6451854ab7e4b4c146bf78a.
% These are local numerical estimates; active rays and sectors remain separate.

\begin{lemma}[A common quarter mesh in a closed neighborhood]
    \label{thm:area_law_near_mark_quarter_mesh}
    \lean{TNLean.PEPS.AreaLaw.Geometry.fineLayer_near_mark_quarter_mesh}
    \leanok
    \uses{def:area_law_fine_layer_indices,def:area_law_belt_cell_marks,
        def:area_law_affine_mesh}
    Let $k\ge 50{,}000{,}000$ and let $z,w$ be fine-cell indices
    in layers $k,h$, respectively. Suppose $v\in M_{k,z}$ and the closed
    cell indexed by $w$ meets the closed square of sup radius $10t$
    about $v$, where $t=2^{\ell_k}$. Then
    \begin{align}
        h                   & \le k+1,                                  & k & \le h+1, \\
        \ell_h              & =\ell_k\quad\text{or}\quad\ell_h+1=\ell_k
        \quad\text{or}\quad\ell_h=\ell_k+1,                                            \\
        M_{k,z}\cup M_{h,w} & \subseteq o+(t/4)\mathbb Z^2.
    \end{align}
\end{lemma}
\begin{proof}\leanok
    \uses{thm:area_law_local_layer_exclusion,thm:area_law_fine_cell_containment,
        thm:area_law_adjacent_fine_meshes,thm:area_law_belt_marks_quarter_mesh}
    The reference mark and a contact point belong to their respective
    closed layers, and their distance is at most $10t$.
    Nonadjacent layers are separated by at least $16t$, so their
    indices differ by at most one. Monotonicity of the fine exponent
    and its successor increment give the three alternatives, so both
    exponents are at least $\ell_k-1$.
    The quarter-mesh containment of belt marks and the nesting of
    adjacent fine meshes therefore place both
    nine-point sets in the same quarter mesh.
\end{proof}

\begin{lemma}[Fine-scale separation from the dummy neighborhood]
    \label{thm:area_law_dummy_fine_scale_separation}
    \lean{TNLean.PEPS.AreaLaw.Geometry.dyadicNeighborhood_dist_later_layer_fineScale}
    \leanok
    \uses{def:area_law_dyadic_layers}
    Suppose $k\ge 50{,}000{,}000$ and $k_0+1\le k$. If
    $x\in\overline{N_{k_0}}$ and $y\in\overline{D_k}$, then
    \begin{align}
        \|x-y\|_\infty & \ge 16\,2^{\ell_k}.
    \end{align}
\end{lemma}
\begin{proof}\leanok
    \uses{thm:area_law_neighborhood_endpoint_distance,thm:area_law_dyadic_layer_distance,
        thm:area_law_scale_side_ratios}
    There is an endpoint $u\in Z$ with
    $\|x-u\|_\infty\le(C+1)2^{k_0}$, while
    $\|y-u\|_\infty\ge C2^k$. The triangle inequality and
    $2\,2^{k_0}\le2^k$ give
    \begin{align}
        2\|x-y\|_\infty & \ge(C-1)2^k\ge2^k.
    \end{align}
    The scale-ratio inequality gives
    $32\,2^{\ell_k}\le2^k$, proving the result.
\end{proof}

